\chapter{Chapter06 — Deterministic Goldreich Levin}
%%% generated by the blueprint pipeline from FABL.Chapter06.LearningAndTesting.DeterministicGoldreichLevin
\section{Overview}

Declarations of the Lean module \texttt{FABL.Chapter06.LearningAndTesting.DeterministicGoldreichLevin}.

\section{Declarations}

\begin{definition}[Deterministic Goldreich–Levin input]
\lean{FABL.DeterministicGoldreichLevinInput}
\statementsource{Analysis-of-Boolean-Functions-by-Ryan-ODonnell}{
  pdf-ffbe9c7b3306-p083-b002,
  pdf-ffbe9c7b3306-p084-b006
}
\label{FABL.DeterministicGoldreichLevinInput}
\leanok
\uses{FABL.SmallBiasInput}
A deterministic Goldreich–Levin input consists of a small-bias input and a
natural-number Fourier $1$-norm bound that is at least $1$.
\end{definition}

\begin{definition}[Rational learning parameter]
\lean{FABL.DeterministicGoldreichLevinInput.learningParameter}
\label{FABL.DeterministicGoldreichLevinInput.learningParameter}
\leanok
\uses{FABL.DeterministicGoldreichLevinInput, FABL.PositiveLearningParameter}
Let the input specify a positive dimension $n$, positive integers $a$ and $b$ with
$2a\le b$, and an integer Fourier bound $B\ge1$. Its learning parameter is the rational
number $a/b$, which lies in $(0,1/2]$.
\end{definition}

\begin{theorem}[Learning accuracy as a real number]
\lean{FABL.DeterministicGoldreichLevinInput.learningParameter_cast}
\label{FABL.DeterministicGoldreichLevinInput.learningParameter_cast}
\leanok
\uses{FABL.DeterministicGoldreichLevinInput, FABL.DeterministicGoldreichLevinInput.learningParameter, FABL.SmallBiasInput.epsilon}
\difficulty{1}
Let a deterministic Goldreich–Levin input have dimension $n>0$, accuracy $a/b$ with
integers $a>0$ and $2a\le b$, and an integer Fourier bound $S\ge1$. Its learning
parameter, viewed as a real number, equals the real accuracy $a/b$.
\end{theorem}

\begin{definition}[Family size bound]
\lean{FABL.DeterministicGoldreichLevinInput.familySizeBound}
\statementsource{Analysis-of-Boolean-Functions-by-Ryan-ODonnell}{
  pdf-ffbe9c7b3306-p083-b002,
  pdf-ffbe9c7b3306-p085-b003
}
\label{FABL.DeterministicGoldreichLevinInput.familySizeBound}
\leanok
\uses{FABL.DeterministicGoldreichLevinInput, FABL.DeterministicGoldreichLevinInput.learningParameter, FABL.fourierOneNormFamilySizeBound}
For a deterministic Goldreich–Levin input with Fourier bound $s$ and learning parameter
determined by its small-bias accuracy, its family-size bound is the value of the
prescribed Fourier $1$-norm family-size bound at $s$ and that learning parameter.
\end{definition}

\begin{theorem}[Positivity of the family size bound]
\lean{FABL.DeterministicGoldreichLevinInput.familySizeBound_pos}
\label{FABL.DeterministicGoldreichLevinInput.familySizeBound_pos}
\leanok
\uses{FABL.DeterministicGoldreichLevinInput, FABL.DeterministicGoldreichLevinInput.familySizeBound, FABL.fourierOneNormFamilySizeBound_pos}
\difficulty{1}
Let a deterministic Goldreich–Levin input have dimension $n>0$, accuracy $a/b$ with
integers $a>0$ and $2a\le b$, and an integer Fourier bound $S\ge1$. For
$\varepsilon=a/b$, the integer $\max\!\left(1,\left\lceil
4S^2/\varepsilon\right\rceil\right)$ is positive.
\end{theorem}

\begin{definition}[Goldreich–Levin threshold]
\lean{FABL.DeterministicGoldreichLevinInput.threshold}
\statementsource{Analysis-of-Boolean-Functions-by-Ryan-ODonnell}{pdf-ffbe9c7b3306-p083-b002}
\label{FABL.DeterministicGoldreichLevinInput.threshold}
\leanok
\uses{FABL.DeterministicGoldreichLevinInput, FABL.DeterministicGoldreichLevinInput.familySizeBound, FABL.DeterministicGoldreichLevinInput.familySizeBound_pos, FABL.DeterministicGoldreichLevinInput.learningParameter, FABL.GoldreichLevinThreshold, FABL.kushilevitzMansourThreshold}
For a deterministic Goldreich–Levin input, its threshold is the Kushilevitz–Mansour
threshold associated with its learning parameter and family-size bound; this threshold
lies in $(0,1]$ and is supplied with the positivity condition on the family-size bound.
\end{definition}

\begin{definition}[Restricted-weight bias parameters]
\lean{FABL.DeterministicGoldreichLevinInput.restrictedWeightBias}
\statementsource{Analysis-of-Boolean-Functions-by-Ryan-ODonnell}{pdf-ffbe9c7b3306-p084-b002}
\label{FABL.DeterministicGoldreichLevinInput.restrictedWeightBias}
\leanok
\uses{FABL.DeterministicGoldreichLevinInput, FABL.DeterministicGoldreichLevinInput.familySizeBound, FABL.DeterministicGoldreichLevinInput.familySizeBound_pos, FABL.SmallBiasInput}
Given a deterministic Goldreich–Levin input with small-bias parameters $(n,a,b)$ and
family-size bound $M$, the restricted-weight bias parameters have dimension $n$,
numerator $a^2$, and denominator $8(4Mb)^2$. Thus the dimension and numerator are
positive and twice the numerator is at most the denominator.
\end{definition}

\begin{definition}[Restricted-weight input]
\lean{FABL.DeterministicGoldreichLevinInput.restrictedWeightInput}
\statementsource{Analysis-of-Boolean-Functions-by-Ryan-ODonnell}{
  pdf-ffbe9c7b3306-p084-b002,
  pdf-ffbe9c7b3306-p084-b004
}
\label{FABL.DeterministicGoldreichLevinInput.restrictedWeightInput}
\leanok
\uses{FABL.DeterministicGoldreichLevinInput, FABL.DeterministicGoldreichLevinInput.restrictedWeightBias, FABL.RestrictedWeightInput}
For a deterministic Goldreich–Levin input, the associated restricted-weight input uses
its restricted-weight bias parameters and the same Fourier $1$-norm bound, which is at
least one.
\end{definition}

\begin{theorem}[Restricted weight accuracy]
\lean{FABL.DeterministicGoldreichLevinInput.restrictedWeightInput_epsilon}
\label{FABL.DeterministicGoldreichLevinInput.restrictedWeightInput_epsilon}
\leanok
\uses{FABL.DeterministicGoldreichLevinInput, FABL.DeterministicGoldreichLevinInput.familySizeBound, FABL.DeterministicGoldreichLevinInput.familySizeBound_pos, FABL.DeterministicGoldreichLevinInput.learningParameter, FABL.DeterministicGoldreichLevinInput.restrictedWeightBias, FABL.DeterministicGoldreichLevinInput.restrictedWeightInput, FABL.DeterministicGoldreichLevinInput.threshold, FABL.RestrictedWeightInput.epsilon, FABL.SmallBiasInput.epsilon, FABL.kushilevitzMansourThreshold}
\difficulty{4}
Let a deterministic Goldreich–Levin input have dimension $n>0$, accuracy $a/b$ with
integers $a>0$ and $2a\le b$, and an integer Fourier bound $S\ge1$. Put
$\varepsilon=a/b$, $M=\max\!\left(1,\left\lceil 4S^2/\varepsilon\right\rceil\right)$,
and $\tau=\varepsilon/(4M)$. The restricted weight input has accuracy $\tau^2/8$.
\end{theorem}

\begin{definition}[Coefficient bias parameters]
\lean{FABL.DeterministicGoldreichLevinInput.coefficientBias}
\statementsource{Analysis-of-Boolean-Functions-by-Ryan-ODonnell}{pdf-ffbe9c7b3306-p083-b002}
\label{FABL.DeterministicGoldreichLevinInput.coefficientBias}
\leanok
\uses{FABL.DeterministicGoldreichLevinInput, FABL.DeterministicGoldreichLevinInput.threshold, FABL.SmallBiasInput, FABL.goldreichLevinActiveCap, FABL.goldreichLevinActiveCap_pos}
For a deterministic Goldreich–Levin input with small-bias parameters $(n,a,b)$ and
threshold $\tau$, the coefficient-bias parameters have dimension $n$, numerator $a$, and
denominator $b$ times twice the active coefficient cap associated with $\tau$. The
dimension and numerator are positive, and twice the numerator is at most the
denominator.
\end{definition}

\begin{definition}[Coefficient-estimation input]
\lean{FABL.DeterministicGoldreichLevinInput.coefficientInput}
\statementsource{Analysis-of-Boolean-Functions-by-Ryan-ODonnell}{
  pdf-ffbe9c7b3306-p080-b002,
  pdf-ffbe9c7b3306-p083-b002
}
\label{FABL.DeterministicGoldreichLevinInput.coefficientInput}
\leanok
\uses{FABL.DeterministicGoldreichLevinInput, FABL.DeterministicGoldreichLevinInput.coefficientBias, FABL.SmallBiasFourierInput}
For a deterministic Goldreich–Levin input, the coefficient-estimation input consists of
its coefficient-bias parameters and its original Fourier $1$-norm bound, which is at
least one.
\end{definition}

\begin{theorem}[Coefficient bias accuracy]
\lean{FABL.DeterministicGoldreichLevinInput.coefficientInput_epsilon}
\label{FABL.DeterministicGoldreichLevinInput.coefficientInput_epsilon}
\leanok
\uses{FABL.DeterministicGoldreichLevinInput, FABL.DeterministicGoldreichLevinInput.coefficientInput, FABL.DeterministicGoldreichLevinInput.learningParameter, FABL.DeterministicGoldreichLevinInput.threshold, FABL.SmallBiasFourierInput.epsilon, FABL.SmallBiasInput.epsilon, FABL.goldreichLevinActiveCap, FABL.goldreichLevinActiveCap_pos}
\difficulty{4}
Let a deterministic Goldreich–Levin input have dimension $n>0$, accuracy $a/b$ with
integers $a>0$ and $2a\le b$, and an integer Fourier bound $S\ge1$. Write
$\varepsilon=a/b$, let $M=\max\!\left(1,\left\lceil
4S^2/\varepsilon\right\rceil\right)$, and set $\tau=\varepsilon/(4M)$. The coefficient
input has accuracy $\varepsilon/(2\lceil 4/\tau^2\rceil)$.
\end{theorem}

\begin{definition}[Rational indexed monomial]
\lean{FABL.rationalIndexedMonomial}
\statementsource{Analysis-of-Boolean-Functions-by-Ryan-ODonnell}{
  pdf-ffbe9c7b3306-p023-b001,
  pdf-ffbe9c7b3306-p083-b001
}
\label{FABL.rationalIndexedMonomial}
\leanok
\uses{FABL.IndexedSignCube, FABL.Sign}
Let $\iota$ be an index set, let $S$ be a finite subset of $\iota$, and let $x$ assign a
sign in $\{ -1,1\}$ to each index. The rational indexed monomial of $x$ on $S$ is
$\prod_{i\in S} x_i$, viewed as a rational number.
\end{definition}

\begin{theorem}[Rational and real indexed monomials]
\lean{FABL.rationalIndexedMonomial_cast}
\label{FABL.rationalIndexedMonomial_cast}
\leanok
\uses{FABL.IndexedSignCube, FABL.indexedMonomial, FABL.rationalIndexedMonomial, FABL.signValue}
\difficulty{1}
Let $\iota$ be any index type, let $S$ be a finite subset of $\iota$, and let
$x:\iota\to\{\pm1\}$ be a sign assignment. The rational product of the signs $x(i)$ over
$i\in S$, viewed as a real number, equals the real product of those signs.
\end{theorem}

\begin{definition}[Restricted Fourier weight estimate from answers]
\lean{FABL.rationalRestrictedFourierWeightEstimateFromAnswers}
\statementsource{Analysis-of-Boolean-Functions-by-Ryan-ODonnell}{
  pdf-ffbe9c7b3306-p084-b002,
  pdf-ffbe9c7b3306-p084-b005
}
\label{FABL.rationalRestrictedFourierWeightEstimateFromAnswers}
\leanok
\uses{FABL.RestrictedWeightInput, FABL.RestrictedWeightInput.innerInput, FABL.RestrictedWeightInput.outerInput, FABL.Sign, FABL.SmallBiasFourierInput.sampleCount, FABL.deterministicRestrictedWeightQueryCount, FABL.rationalIndexedMonomial, FABL.restrictedWeightFreeAssignment}
Let the input consist of a dimension $n>0$, positive integers $a,b$ with $2a\le b$, and
an integer Fourier bound $B\ge1$. For a coordinate set $J\subseteq\mathrm{Fin}(n)$, a
subset $S\subseteq J$, and a sign-valued answer for each of the prescribed outer–inner
query pairs, define the estimate as follows: for each outer sample, average the answers
multiplied by the sign monomial indexed by $S$ on the corresponding inner free
assignment; square this average, then average the squares over all outer samples.
\end{definition}

\begin{theorem}[Casting the restricted weight estimate]
\lean{FABL.rationalRestrictedFourierWeightEstimateFromAnswers_cast}
\label{FABL.rationalRestrictedFourierWeightEstimateFromAnswers_cast}
\leanok
\uses{FABL.RestrictedWeightInput, FABL.Sign, FABL.deterministicRestrictedWeightEstimateFromAnswers, FABL.deterministicRestrictedWeightQueryCount, FABL.rationalRestrictedFourierWeightEstimateFromAnswers, FABL.signValue}
\difficulty{1}
Let $n,a,b,s$ be natural numbers with $0<n$, $0<a$, $2a\le b$, and $1\le s$. Consider
the restricted-weight input specified by dimension $n$, bias numerator $a$, bias
denominator $b$, and Fourier bound $s$. For any coordinate set
$J\subseteq\{1,\ldots,n\}$, any $S\subseteq J$, and any assignment of signs to the
queries in the associated finite query set, the rational estimate obtained by averaging
the squared inner sample averages over the outer samples, divided by the outer sample
count, equals the corresponding real-valued estimate.
\end{theorem}

\begin{theorem}[Casting the oracle restricted weight estimate]
\lean{FABL.rationalRestrictedFourierWeightEstimateFromOracle_cast}
\label{FABL.rationalRestrictedFourierWeightEstimateFromOracle_cast}
\leanok
\uses{FABL.BooleanFunction, FABL.BooleanFunction.toReal, FABL.RestrictedWeightInput, FABL.deterministicRestrictedWeightEstimateFromAnswers, FABL.deterministicRestrictedWeightQueryPair, FABL.deterministicRestrictedWeightQueryPoint, FABL.rationalRestrictedFourierWeightEstimateFromAnswers, FABL.rationalRestrictedFourierWeightEstimateFromAnswers_cast, FABL.restrictedFourierCoefficientEstimate, FABL.restrictedFourierWeightEstimate}
\difficulty{3}
Let $n,a,b,s$ be natural numbers with $0<n$, $0<a$, $2a\le b$, and $1\le s$. Consider
the restricted-weight input specified by dimension $n$, bias numerator $a$, bias
denominator $b$, and Fourier bound $s$, and let $f$ be a sign-valued function on the
$n$-dimensional sign cube. For any coordinate set $J\subseteq\{1,\ldots,n\}$ and any
$S\subseteq J$, the rational restricted weight estimate obtained by querying $f$ at
every point of the prescribed query set, when viewed as a real number, equals the
corresponding real restricted Fourier weight estimate of $f$.
\end{theorem}

\begin{definition}[Rational small-bias Fourier estimate]
\lean{FABL.rationalSmallBiasFourierEstimate}
\statementsource{Analysis-of-Boolean-Functions-by-Ryan-ODonnell}{
  pdf-ffbe9c7b3306-p080-b002,
  pdf-ffbe9c7b3306-p080-b003
}
\label{FABL.rationalSmallBiasFourierEstimate}
\leanok
\uses{FABL.Sign, FABL.SmallBiasFourierInput, FABL.SmallBiasFourierInput.generatorInput, FABL.SmallBiasFourierInput.sample, FABL.SmallBiasFourierInput.sampleCount, FABL.binaryCubeSignEquiv, FABL.empiricalFourierCoeff}
Let the input consist of a dimension $n>0$, positive integers $a,b$ with $2a\le b$, and
an integer Fourier bound $B\ge1$. For each coordinate set $U\subseteq\mathrm{Fin}(n)$
and each sign-valued answer indexed by the sample set, define the estimate to be the
empirical Fourier coefficient on $U$ of the paired samples: the signs obtained by
encoding the generated binary vectors, paired with their answers.
\end{definition}

\begin{theorem}[Casting the small-bias Fourier estimate]
\lean{FABL.rationalSmallBiasFourierEstimate_cast}
\label{FABL.rationalSmallBiasFourierEstimate_cast}
\leanok
\uses{FABL.BooleanFunction, FABL.BooleanFunction.toReal, FABL.SmallBiasFourierInput, FABL.SmallBiasFourierInput.generatorInput, FABL.SmallBiasFourierInput.sample, FABL.binaryCubeSignEquiv, FABL.empiricalFourierCoeff_cast, FABL.fourierObservation, FABL.rationalSmallBiasFourierEstimate, FABL.realEmpiricalFourierCoeff, FABL.smallBiasFourierEstimate}
\difficulty{3}
Let $n,a,b,s$ be natural numbers with $0<n$, $0<a$, $2a\le b$, and $1\le s$, and
consider the small-bias Fourier input with dimension $n$, bias numerator $a$, bias
denominator $b$, and Fourier bound $s$. Let $f$ be a sign-valued function on the
$n$-dimensional sign cube and let $U$ be a set of coordinates in the generator's output
dimension. The rational empirical Fourier coefficient obtained from the generated
samples labeled by $f$, when viewed as a real number, equals the real sample Fourier
estimate of $f$ on $U$.
\end{theorem}

\begin{definition}[Goldreich–Levin stage records from answers]
\lean{FABL.deterministicGoldreichLevinStageRecordsFromAnswers}
\statementsource{Analysis-of-Boolean-Functions-by-Ryan-ODonnell}{
  pdf-ffbe9c7b3306-p085-b001,
  pdf-ffbe9c7b3306-p085-b004
}
\label{FABL.deterministicGoldreichLevinStageRecordsFromAnswers}
\leanok
\uses{FABL.DeterministicGoldreichLevinInput, FABL.DeterministicGoldreichLevinInput.restrictedWeightInput, FABL.GoldreichLevinEstimateRecord, FABL.Sign, FABL.deterministicRestrictedWeightQueryCount, FABL.freeFrequencyPart, FABL.goldreichLevinCandidates, FABL.prefixCoordinates, FABL.rationalRestrictedFourierWeightEstimateFromAnswers}
Let the input consist of a dimension $n>0$, positive accuracy numerator and denominator
$a,b$ with $2a\le b$, and an integer Fourier bound $B\ge1$. Given an index $k<n$, a
finite family of coordinate subsets of $\mathrm{Fin}(n)$, and one sign-valued answer for
each restricted-weight query, form a list with one record for each candidate obtained
from the family by either retaining a set or adjoining coordinate $k$. Each record uses
the prefix of coordinates $0,\ldots,k$ and stores, as its frequency prefix, that
candidate; its rational estimate is the restricted Fourier weight estimate for the
candidate's part lying in the prefix.
\end{definition}

\begin{definition}[Goldreich–Levin stage records for a target]
\lean{FABL.deterministicGoldreichLevinStageRecords}
\statementsource{Analysis-of-Boolean-Functions-by-Ryan-ODonnell}{
  pdf-ffbe9c7b3306-p085-b001,
  pdf-ffbe9c7b3306-p085-b004
}
\label{FABL.deterministicGoldreichLevinStageRecords}
\leanok
\uses{FABL.BooleanFunction, FABL.DeterministicGoldreichLevinInput, FABL.DeterministicGoldreichLevinInput.restrictedWeightInput, FABL.GoldreichLevinEstimateRecord, FABL.deterministicGoldreichLevinStageRecordsFromAnswers, FABL.deterministicRestrictedWeightQueryPoint, FABL.prefixCoordinates}
Let the input consist of a dimension $n>0$, positive accuracy numerator and denominator
$a,b$ with $2a\le b$, and an integer Fourier bound $B\ge1$, and let
$f:\{−1,1\}^n\to\{−1,1\}$ be a Boolean function. Given $k<n$ and a finite active family
of coordinate subsets of $\mathrm{Fin}(n)$, form the stage records as in the
answer-based construction, using as each query answer the value of $f$ at the
corresponding restricted-weight query point.
\end{definition}

\begin{theorem}[Accuracy of Goldreich–Levin stage records]
\lean{FABL.deterministicGoldreichLevinStageRecords_isAccurate}
\proofsource{Analysis-of-Boolean-Functions-by-Ryan-ODonnell}{
  pdf-ffbe9c7b3306-p085-b002,
  pdf-ffbe9c7b3306-p085-b004
}
\label{FABL.deterministicGoldreichLevinStageRecords_isAccurate}
\leanok
\uses{FABL.BooleanFunction, FABL.BooleanFunction.toReal, FABL.DeterministicGoldreichLevinInput, FABL.DeterministicGoldreichLevinInput.restrictedWeightInput, FABL.DeterministicGoldreichLevinInput.restrictedWeightInput_epsilon, FABL.DeterministicGoldreichLevinInput.threshold, FABL.GoldreichLevinEstimateRecord.IsAccurate, FABL.GoldreichLevinEstimateTrace.IsAccurate, FABL.RestrictedWeightInput.epsilon, FABL.abs_restrictedFourierWeightEstimate_sub_le, FABL.deterministicGoldreichLevinStageRecords, FABL.deterministicGoldreichLevinStageRecordsFromAnswers, FABL.deterministicRestrictedWeightQueryPoint, FABL.fourierOneNorm, FABL.freeFrequencyPart, FABL.prefixCoordinates, FABL.rationalRestrictedFourierWeightEstimateFromAnswers, FABL.rationalRestrictedFourierWeightEstimateFromOracle_cast, FABL.restrictedFourierWeight, FABL.restrictedFourierWeightEstimate}
\difficulty{4}
Let the input specify an output dimension $n>0$, a positive rational accuracy at most
$1/2$, and an integral Fourier $1$-norm bound $s\ge1$. Let $f$ be a sign-valued function
on the $n$-dimensional sign cube whose Fourier $1$-norm is at most $s$. For every $k<n$
and every finite family $A$ of coordinate subsets, each estimate record produced at
stage $k$ differs from the corresponding restricted Fourier weight of $f$ by less than
one quarter of the square of the learning threshold specified by the input.
\end{theorem}

\begin{theorem}[Coverage of Goldreich–Levin stage records]
\lean{FABL.deterministicGoldreichLevinStageRecords_covers}
\label{FABL.deterministicGoldreichLevinStageRecords_covers}
\leanok
\uses{FABL.BooleanFunction, FABL.DeterministicGoldreichLevinInput, FABL.DeterministicGoldreichLevinInput.restrictedWeightInput, FABL.GoldreichLevinEstimateRecord.Matches, FABL.GoldreichLevinEstimateTrace.Covers, FABL.deterministicGoldreichLevinStageRecords, FABL.deterministicGoldreichLevinStageRecordsFromAnswers, FABL.deterministicRestrictedWeightQueryPoint, FABL.freeFrequencyPart, FABL.goldreichLevinCandidates, FABL.prefixCoordinates, FABL.rationalRestrictedFourierWeightEstimateFromAnswers}
\difficulty{4}
Let the input specify an output dimension $n>0$, a positive rational accuracy at most
$1/2$, and an integral Fourier $1$-norm bound at least $1$. Let $f$ be a sign-valued
function on the $n$-dimensional sign cube, let $k<n$, and let $A$ be a family of subsets
of the $n$ coordinates. At stage $k$, every candidate obtained from a member of $A$ by
either leaving it unchanged or adjoining coordinate $k$ occurs as a frequency prefix in
some estimate record whose coordinates are precisely the first $k+1$ coordinates.
\end{theorem}

\begin{definition}[Deterministic Goldreich–Levin state sequence]
\lean{FABL.deterministicGoldreichLevinState}
\statementsource{Analysis-of-Boolean-Functions-by-Ryan-ODonnell}{
  pdf-ffbe9c7b3306-p085-b001,
  pdf-ffbe9c7b3306-p085-b004
}
\label{FABL.deterministicGoldreichLevinState}
\leanok
\uses{FABL.BooleanFunction, FABL.DeterministicGoldreichLevinInput, FABL.DeterministicGoldreichLevinInput.threshold, FABL.GoldreichLevinQueryState, FABL.deterministicGoldreichLevinStageRecords, FABL.goldreichLevinActiveCap, FABL.goldreichLevinQueryStep}
Let $n>0$, and let $a,b$ be positive integers with $2a\le b$, together with an integer
Fourier bound $B\ge1$. For a Boolean function $f:\{−1,1\}^n\to\{−1,1\}$, define a
sequence of states indexed by the nonnegative integers. The initial state has active
family $\{\varnothing\}$ and empty trace; at stage $k+1$, if $k<n$ and the preceding
active family exists with size at most the active cap determined by the input threshold,
update the state by the Goldreich–Levin query step at coordinate $k$, using the estimate
records computed from $f$; if that family exceeds the cap, make the active family absent
and preserve the trace. In all other cases, retain the preceding state unchanged.
\end{definition}

\begin{theorem}[Deterministic goldreich levin state spec]
\lean{FABL.deterministicGoldreichLevinState_spec}
\statementsource{Analysis-of-Boolean-Functions-by-Ryan-ODonnell}{
  pdf-ffbe9c7b3306-p085-b001,
  pdf-ffbe9c7b3306-p085-b004
}
\label{FABL.deterministicGoldreichLevinState_spec}
\leanok
\uses{FABL.BooleanFunction, FABL.BooleanFunction.toReal, FABL.DeterministicGoldreichLevinInput, FABL.DeterministicGoldreichLevinInput.threshold, FABL.GoldreichLevinEstimateTrace.Covers, FABL.GoldreichLevinEstimateTrace.IsAccurate, FABL.GoldreichLevinQueryState.SatisfiesInvariant, FABL.deterministicGoldreichLevinStageRecords, FABL.deterministicGoldreichLevinStageRecords_covers, FABL.deterministicGoldreichLevinStageRecords_isAccurate, FABL.deterministicGoldreichLevinState, FABL.fourierOneNorm, FABL.goldreichLevinCandidates, FABL.goldreichLevinInitialQueryState_satisfiesInvariant, FABL.goldreichLevinQueryStep, FABL.goldreichLevinQueryStep_satisfiesInvariant, FABL.prefixCoordinates}
\difficulty{5}
% dataset-pipeline: stub -- prose pending build_informal_statements.py
Statement pending regeneration.
\end{theorem}

\begin{theorem}[Deterministic goldreich levin state is correct output]
\lean{FABL.deterministicGoldreichLevinState_isCorrectOutput}
\statementsource{Analysis-of-Boolean-Functions-by-Ryan-ODonnell}{
  pdf-ffbe9c7b3306-p083-b002,
  pdf-ffbe9c7b3306-p085-b002
}
\label{FABL.deterministicGoldreichLevinState_isCorrectOutput}
\leanok
\uses{FABL.BooleanFunction, FABL.BooleanFunction.toReal, FABL.DeterministicGoldreichLevinInput, FABL.DeterministicGoldreichLevinInput.threshold, FABL.GoldreichLevinQueryState.IsCorrectOutput, FABL.deterministicGoldreichLevinState, FABL.deterministicGoldreichLevinState_spec, FABL.fourierOneNorm, FABL.goldreichLevinActiveInvariant_card_le, FABL.goldreichLevinActiveInvariant_complete, FABL.goldreichLevinActiveInvariant_sound}
\difficulty{3}
% dataset-pipeline: stub -- prose pending build_informal_statements.py
Statement pending regeneration.
\end{theorem}

\begin{definition}[Restricted local work]
\lean{FABL.deterministicGoldreichLevinRestrictedLocalWork}
\label{FABL.deterministicGoldreichLevinRestrictedLocalWork}
\leanok
\uses{FABL.DeterministicGoldreichLevinInput, FABL.DeterministicGoldreichLevinInput.restrictedWeightInput, FABL.RestrictedWeightInput.innerCount, FABL.RestrictedWeightInput.outerCount}
For an input specifying a positive dimension, positive rational bias at most one half,
and an integral Fourier bound at least one, let $O$ and $I$ denote the outer and inner
counts in its restricted-weight data. The restricted local-work quantity is
$O\bigl(I(n+2)+3\bigr)+1$, where $n$ is the input dimension.
\end{definition}

\begin{definition}[Stage local work]
\lean{FABL.deterministicGoldreichLevinStageLocalWork}
\label{FABL.deterministicGoldreichLevinStageLocalWork}
\leanok
\uses{FABL.DeterministicGoldreichLevinInput, FABL.deterministicGoldreichLevinRestrictedLocalWork, FABL.goldreichLevinCandidates}
For an input specifying a positive dimension, positive rational bias at most one half,
and an integral Fourier bound at least one, let $k<n$ and let $A$ be a family of subsets
of the $n$ coordinates. The stage local-work quantity is the number of Goldreich–Levin
candidates at stage $k$ determined by $A$, multiplied by the input’s restricted
local-work quantity.
\end{definition}

\begin{definition}[Deterministic stage cost]
\lean{FABL.deterministicGoldreichLevinStageCost}
\label{FABL.deterministicGoldreichLevinStageCost}
\leanok
\uses{FABL.DeterministicGoldreichLevinInput, FABL.DeterministicGoldreichLevinInput.restrictedWeightInput, FABL.DeterministicQueryProgram.queryBatchCost, FABL.LearningCost, FABL.deterministicGoldreichLevinStageLocalWork, FABL.deterministicRestrictedWeightQueryCount}
Let the output dimension be a positive integer $n$, let the requested bias be a positive
rational number at most $1/2$, and let the integral Fourier $1$-norm bound be an integer
at least $1$. For an integer $k<n$ and a finite family of subsets of the $n$
coordinates, define the stage cost as the restricted-weight query-batch cost plus a cost
with zero random examples and zero membership queries that charges the stage’s local
computation in work.
\end{definition}

\begin{definition}[Goldreich–Levin stage query program]
\lean{FABL.deterministicGoldreichLevinStageProgram}
\statementsource{Analysis-of-Boolean-Functions-by-Ryan-ODonnell}{
  pdf-ffbe9c7b3306-p085-b001,
  pdf-ffbe9c7b3306-p085-b004
}
\label{FABL.deterministicGoldreichLevinStageProgram}
\leanok
\uses{FABL.DeterministicGoldreichLevinInput, FABL.DeterministicGoldreichLevinInput.restrictedWeightInput, FABL.DeterministicQueryProgram, FABL.GoldreichLevinEstimateRecord, FABL.Sign, FABL.deterministicGoldreichLevinStageLocalWork, FABL.deterministicGoldreichLevinStageRecordsFromAnswers, FABL.deterministicRestrictedWeightQueryCount, FABL.deterministicRestrictedWeightQueryPoint, FABL.prefixCoordinates}
Let the input consist of a dimension $n>0$, positive accuracy numerator and denominator
$a,b$ with $2a\le b$, and an integer Fourier bound $B\ge1$. For $k<n$ and a finite
active family of coordinate subsets, the stage program queries the full batch of
restricted-weight query points for the prefix $0,\ldots,k$. Once the answers are
supplied, it charges the prescribed local computation cost for that stage and returns
the corresponding list of estimate records.
\end{definition}

\begin{theorem}[Cost of a Goldreich–Levin stage]
\lean{FABL.DeterministicQueryProgram.runWithCost_deterministicGoldreichLevinStageProgram}
\label{FABL.DeterministicQueryProgram.runWithCost_deterministicGoldreichLevinStageProgram}
\leanok
\uses{FABL.BooleanFunction, FABL.DeterministicGoldreichLevinInput, FABL.DeterministicQueryProgram.runWithCost, FABL.deterministicGoldreichLevinStageCost, FABL.deterministicGoldreichLevinStageProgram, FABL.deterministicGoldreichLevinStageRecords}
\difficulty{1}
Let $n$ be a positive integer, let $a,b$ be positive integers with $2a\le b$, and let
$M\ge1$ be an integer Fourier norm bound. For a Boolean target
$f:\{\pm1\}^n\to\{\pm1\}$, a stage index $k<n$, and a finite active family of coordinate
subsets, running the deterministic Goldreich–Levin stage with oracle $f$ returns its
prescribed list of estimate records and has exactly the prescribed stage cost: one batch
of restricted-weight queries together with the stage’s charged local work.
\end{theorem}

\begin{definition}[Goldreich–Levin stage step program]
\lean{FABL.deterministicGoldreichLevinStageStepProgram}
\statementsource{Analysis-of-Boolean-Functions-by-Ryan-ODonnell}{pdf-ffbe9c7b3306-p085-b001}
\label{FABL.deterministicGoldreichLevinStageStepProgram}
\leanok
\uses{FABL.DeterministicGoldreichLevinInput, FABL.DeterministicGoldreichLevinInput.threshold, FABL.DeterministicQueryProgram, FABL.DeterministicQueryProgram.map, FABL.GoldreichLevinQueryState, FABL.Sign, FABL.deterministicGoldreichLevinStageProgram, FABL.goldreichLevinControllerStepWork, FABL.goldreichLevinQueryStep}
Let the input consist of a dimension $n>0$, positive accuracy numerator and denominator
$a,b$ with $2a\le b$, and an integer Fourier bound $B\ge1$. Given $k<n$, a current query
state, and a finite active family of coordinate subsets, the step program first charges
the prescribed controller computation cost, then runs the stage query program and maps
its output records to the state obtained by applying the Goldreich–Levin query step to
the current state.
\end{definition}

\begin{theorem}[Output of a Goldreich–Levin stage step]
\lean{FABL.DeterministicQueryProgram.runWithCost_deterministicGoldreichLevinStageStepProgram_fst}
\label{FABL.DeterministicQueryProgram.runWithCost_deterministicGoldreichLevinStageStepProgram_fst}
\leanok
\uses{FABL.BooleanFunction, FABL.DeterministicGoldreichLevinInput, FABL.DeterministicGoldreichLevinInput.threshold, FABL.DeterministicQueryProgram.map, FABL.DeterministicQueryProgram.runWithCost, FABL.DeterministicQueryProgram.runWithCost_bind, FABL.DeterministicQueryProgram.runWithCost_deterministicGoldreichLevinStageProgram, FABL.GoldreichLevinQueryState, FABL.deterministicGoldreichLevinStageRecords, FABL.deterministicGoldreichLevinStageStepProgram, FABL.goldreichLevinQueryStep}
\difficulty{3}
Let $n$ be a positive integer, let $a,b$ be positive integers with $2a\le b$, and let
$M\ge1$ be an integer Fourier norm bound. For a Boolean target
$f:\{\pm1\}^n\to\{\pm1\}$, a stage index $k<n$, a Goldreich–Levin query state, and a
finite active family of coordinate subsets, the output of running the stage-step program
is the state obtained by appending the stage’s estimate records and applying the
prescribed state update.
\end{theorem}

\begin{definition}[Deterministic Goldreich–Levin run-up program]
\lean{FABL.deterministicGoldreichLevinRunUpToProgram}
\statementsource{Analysis-of-Boolean-Functions-by-Ryan-ODonnell}{
  pdf-ffbe9c7b3306-p085-b001,
  pdf-ffbe9c7b3306-p085-b004
}
\label{FABL.deterministicGoldreichLevinRunUpToProgram}
\leanok
\uses{FABL.DeterministicGoldreichLevinInput, FABL.DeterministicGoldreichLevinInput.threshold, FABL.DeterministicQueryProgram, FABL.DeterministicQueryProgram.bind, FABL.GoldreichLevinQueryState, FABL.Sign, FABL.deterministicGoldreichLevinStageStepProgram, FABL.goldreichLevinActiveCap}
Given a deterministic Goldreich–Levin input and a natural number $m$, define a
deterministic query program recursively. At $m=0$, the program returns the state whose
active family is the singleton containing the empty set and whose trace is empty. For
$m=k+1$, it first runs the program for $k$; if $k$ is less than the input dimension and
the resulting state has an active family, then it runs the stage program when that
family has size at most the active cap, and otherwise returns a state with no active
family and the existing trace. If $k$ is outside the input dimension or the resulting
state has no active family, it returns that state unchanged.
\end{definition}

\begin{theorem}[Output of the truncated query run]
\lean{FABL.DeterministicQueryProgram.runWithCost_deterministicGoldreichLevinRunUpToProgram_fst}
\label{FABL.DeterministicQueryProgram.runWithCost_deterministicGoldreichLevinRunUpToProgram_fst}
\leanok
\uses{FABL.BooleanFunction, FABL.DeterministicGoldreichLevinInput, FABL.DeterministicGoldreichLevinInput.threshold, FABL.DeterministicQueryProgram.runWithCost, FABL.DeterministicQueryProgram.runWithCost_bind, FABL.DeterministicQueryProgram.runWithCost_deterministicGoldreichLevinStageStepProgram_fst, FABL.deterministicGoldreichLevinRunUpToProgram, FABL.deterministicGoldreichLevinStageStepProgram, FABL.deterministicGoldreichLevinState, FABL.goldreichLevinActiveCap}
\difficulty{5}
Let $n,a,d,s$ be positive integers with $2a\le d$ and $s\ge1$, and let the target be a
Boolean function on the sign cube of dimension $n$. For every integer $k\ge0$, running
the deterministic Goldreich–Levin query program truncated at stage $k$ returns, as its
output component, the state obtained by the following recursion: start with active
family $\{\varnothing\}$ and empty trace; at each stage below $n$, if the active family
is absent, keep the state, and if its size exceeds the active cap, replace it by the
absent state while retaining the trace. Otherwise append the stage's estimate records to
the trace and refine the active family using those estimates. After stage $n$, keep the
state unchanged.
\end{theorem}

\begin{definition}[Deterministic Goldreich–Levin Program]
\lean{FABL.deterministicGoldreichLevinProgram}
\statementsource{Analysis-of-Boolean-Functions-by-Ryan-ODonnell}{
  pdf-ffbe9c7b3306-p083-b002,
  pdf-ffbe9c7b3306-p085-b001,
  pdf-ffbe9c7b3306-p158-b001,
  pdf-ffbe9c7b3306-p158-b005,
  pdf-ffbe9c7b3306-p158-b006
}
\label{FABL.deterministicGoldreichLevinProgram}
\leanok
\uses{FABL.DeterministicGoldreichLevinInput, FABL.DeterministicGoldreichLevinInput.restrictedWeightInput, FABL.DeterministicQueryProgram, FABL.GoldreichLevinQueryState, FABL.Sign, FABL.deterministicGoldreichLevinRunUpToProgram, FABL.deterministicRestrictedWeightConstructionWork}
Given a deterministic Goldreich–Levin input, define a deterministic query program that
first charges the restricted-weight construction work for that input, then runs the
Goldreich–Levin run-up program for as many steps as the input dimension.
\end{definition}

\begin{theorem}[Deterministic goldreich levin program is correct output]
\lean{FABL.deterministicGoldreichLevinProgram_isCorrectOutput}
\statementsource{Analysis-of-Boolean-Functions-by-Ryan-ODonnell}{
  pdf-ffbe9c7b3306-p083-b002,
  pdf-ffbe9c7b3306-p085-b002,
  pdf-ffbe9c7b3306-p158-b006
}
\label{FABL.deterministicGoldreichLevinProgram_isCorrectOutput}
\leanok
\uses{FABL.BooleanFunction, FABL.BooleanFunction.toReal, FABL.DeterministicGoldreichLevinInput, FABL.DeterministicGoldreichLevinInput.threshold, FABL.DeterministicQueryProgram.runWithCost, FABL.DeterministicQueryProgram.runWithCost_deterministicGoldreichLevinRunUpToProgram_fst, FABL.GoldreichLevinQueryState.IsCorrectOutput, FABL.deterministicGoldreichLevinProgram, FABL.deterministicGoldreichLevinState_isCorrectOutput, FABL.fourierOneNorm}
\difficulty{3}
% dataset-pipeline: stub -- prose pending build_informal_statements.py
Statement pending regeneration.
\end{theorem}

\begin{definition}[Coefficient local computation work]
\lean{FABL.deterministicGoldreichLevinCoefficientLocalWork}
\label{FABL.deterministicGoldreichLevinCoefficientLocalWork}
\leanok
\uses{FABL.DeterministicGoldreichLevinInput, FABL.DeterministicGoldreichLevinInput.coefficientInput, FABL.SmallBiasFourierInput.sampleCount, FABL.smallBiasFourierEstimatorWork}
Let the output dimension be a positive integer $n$, let the requested bias be a positive
rational number at most $1/2$, and let the integral Fourier $1$-norm bound be an integer
at least $1$. Define the coefficient local computation work to be the work of the
small-bias Fourier estimator, using its sample count and the full set of $n$
coordinates.
\end{definition}

\begin{definition}[Estimated Sparse Fourier Hypothesis]
\lean{FABL.deterministicGoldreichLevinHypothesisFromAnswers}
\statementsource{Analysis-of-Boolean-Functions-by-Ryan-ODonnell}{
  pdf-ffbe9c7b3306-p080-b002,
  pdf-ffbe9c7b3306-p080-b003
}
\label{FABL.deterministicGoldreichLevinHypothesisFromAnswers}
\leanok
\uses{FABL.DeterministicGoldreichLevinInput, FABL.DeterministicGoldreichLevinInput.coefficientInput, FABL.Sign, FABL.SmallBiasFourierInput.sampleCount, FABL.SparseFourierHypothesis, FABL.SparseFourierHypothesis.ofCoefficients, FABL.rationalSmallBiasFourierEstimate}
Given a deterministic Goldreich–Levin input, a finite family of frequency sets in the
input dimension, and a sign-valued answer for each sample in the input's
coefficient-estimation sample set, define a sparse Fourier hypothesis supported on that
family, with each coefficient given by the rational small-bias Fourier estimate for the
corresponding frequency set.
\end{definition}

\begin{definition}[Deterministic coefficient cost]
\lean{FABL.deterministicGoldreichLevinCoefficientCost}
\label{FABL.deterministicGoldreichLevinCoefficientCost}
\leanok
\uses{FABL.DeterministicGoldreichLevinInput, FABL.DeterministicGoldreichLevinInput.coefficientInput, FABL.DeterministicQueryProgram.queryBatchCost, FABL.LearningCost, FABL.SmallBiasFourierInput.generatorInput, FABL.SmallBiasFourierInput.sampleCount, FABL.deterministicGoldreichLevinCoefficientLocalWork, FABL.deterministicSmallBiasWork}
Let the output dimension be a positive integer $n$, let the requested bias be a positive
rational number at most $1/2$, and let the integral Fourier $1$-norm bound be an integer
at least $1$. For a finite family of subsets of the $n$ coordinates, define the
coefficient cost as the sum of a cost charging the deterministic small-bias computation
in work and the coefficient query-batch cost, augmented by work equal to the family’s
cardinality times the coefficient local computation work; the added work charges have
zero random examples and zero membership queries.
\end{definition}

\begin{definition}[Goldreich–Levin Coefficient Program]
\lean{FABL.deterministicGoldreichLevinCoefficientProgram}
\statementsource{Analysis-of-Boolean-Functions-by-Ryan-ODonnell}{
  pdf-ffbe9c7b3306-p080-b002,
  pdf-ffbe9c7b3306-p080-b003,
  pdf-ffbe9c7b3306-p157-b008,
  pdf-ffbe9c7b3306-p157-b009
}
\label{FABL.deterministicGoldreichLevinCoefficientProgram}
\leanok
\uses{FABL.DeterministicGoldreichLevinInput, FABL.DeterministicGoldreichLevinInput.coefficientInput, FABL.DeterministicQueryProgram, FABL.Sign, FABL.SmallBiasFourierInput.generatorInput, FABL.SmallBiasFourierInput.sample, FABL.SmallBiasFourierInput.sampleCount, FABL.SparseFourierHypothesis, FABL.binaryCubeSignEquiv, FABL.deterministicGoldreichLevinCoefficientLocalWork, FABL.deterministicGoldreichLevinHypothesisFromAnswers, FABL.deterministicSmallBiasWork}
Given a deterministic Goldreich–Levin input and a finite family of frequency sets in its
dimension, define a deterministic query program that first charges the small-bias
construction work for the coefficient-estimation generator. It then queries the sign
cube at every sample generated by that input, charges local work equal to the family
size times the coefficient-estimation local-work bound, and returns the sparse Fourier
hypothesis formed from the answers.
\end{definition}

\begin{theorem}[Coefficient program output and cost]
\lean{FABL.DeterministicQueryProgram.runWithCost_deterministicGoldreichLevinCoefficientProgram}
\label{FABL.DeterministicQueryProgram.runWithCost_deterministicGoldreichLevinCoefficientProgram}
\leanok
\uses{FABL.BooleanFunction, FABL.DeterministicGoldreichLevinInput, FABL.DeterministicGoldreichLevinInput.coefficientInput, FABL.DeterministicQueryProgram.runWithCost, FABL.SmallBiasFourierInput.sample, FABL.binaryCubeSignEquiv, FABL.deterministicGoldreichLevinCoefficientCost, FABL.deterministicGoldreichLevinCoefficientProgram, FABL.deterministicGoldreichLevinHypothesisFromAnswers}
\difficulty{1}
Let $n,a,d,s$ be positive integers with $2a\le d$ and $s\ge1$, and consider a Boolean
target on the sign cube of dimension $n$, together with a finite family $\mathcal F$ of
subsets of the coordinates. The coefficient program returns the sparse Fourier
hypothesis supported on $\mathcal F$ whose coefficient at each $U\in\mathcal F$ is the
empirical average, over the program's small-bias sample, of the target value times the
sign character indexed by $U$. Its cost has zero random examples, a number of membership
queries equal to the sample count, and local work equal to the sample count plus the
generator construction work and $|\mathcal F|$ times the local coefficient-estimation
work.
\end{theorem}

\begin{theorem}[Coefficient estimation accuracy]
\lean{FABL.deterministicGoldreichLevinHypothesis_coefficients_le}
\label{FABL.deterministicGoldreichLevinHypothesis_coefficients_le}
\leanok
\uses{FABL.BooleanFunction, FABL.BooleanFunction.toReal, FABL.DeterministicGoldreichLevinInput, FABL.DeterministicGoldreichLevinInput.coefficientInput, FABL.DeterministicGoldreichLevinInput.coefficientInput_epsilon, FABL.DeterministicGoldreichLevinInput.learningParameter, FABL.DeterministicGoldreichLevinInput.threshold, FABL.SmallBiasFourierInput.epsilon, FABL.SmallBiasFourierInput.sample, FABL.abs_deterministicSmallBiasFourierEstimate_sub_fourierCoeff_le, FABL.binaryCubeSignEquiv, FABL.finiteFamilyCoefficientAccuracy, FABL.fourierCoeff, FABL.fourierOneNorm, FABL.goldreichLevinActiveCap, FABL.goldreichLevinActiveCap_pos, FABL.rationalSmallBiasFourierEstimate, FABL.rationalSmallBiasFourierEstimate_cast, FABL.smallBiasFourierEstimate}
\difficulty{4}
Let a deterministic Goldreich–Levin input have dimension $n>0$, rational accuracy $a/b$
with $a>0$ and $2a\le b$, and Fourier bound $B\ge1$, and put $\varepsilon=a/b$. Let $f$
be a Boolean function on $n$ sign-valued coordinates with Fourier $1$-norm at most $B$.
For a nonempty finite family $\mathcal F$ of coordinate subsets, assume $|\mathcal
F|\le\lceil4/\tau^2\rceil$, where $M=\max(1,\lceil4B^2/\varepsilon\rceil)$ and
$\tau=\varepsilon/(4M)$. For every $U\in\mathcal F$, the rational sample estimate of the
Fourier coefficient of $f$ at $U$ differs from that coefficient by at most
$\varepsilon/(2|\mathcal F|)$.
\end{theorem}

\begin{definition}[Second-Stage Fourier Hypothesis Program]
\lean{FABL.deterministicGoldreichLevinSecondStage}
\statementsource{Analysis-of-Boolean-Functions-by-Ryan-ODonnell}{
  pdf-ffbe9c7b3306-p080-b006,
  pdf-ffbe9c7b3306-p081-b001
}
\label{FABL.deterministicGoldreichLevinSecondStage}
\leanok
\uses{FABL.DeterministicGoldreichLevinInput, FABL.DeterministicGoldreichLevinInput.threshold, FABL.DeterministicQueryProgram, FABL.GoldreichLevinQueryState, FABL.Sign, FABL.SparseFourierHypothesis, FABL.SparseFourierHypothesis.empty, FABL.deterministicGoldreichLevinCoefficientProgram, FABL.goldreichLevinActiveCap}
Given a deterministic Goldreich–Levin input and a query state, define a deterministic
query program returning a sparse Fourier hypothesis in the input dimension. If the state
has no active family, or its active family has size greater than the active cap at the
input's threshold, the program returns the empty hypothesis; otherwise it runs the
coefficient program on that family.
\end{definition}

\begin{definition}[Deterministic Goldreich–Levin Learner]
\lean{FABL.deterministicGoldreichLevinLearner}
\statementsource{Analysis-of-Boolean-Functions-by-Ryan-ODonnell}{
  pdf-ffbe9c7b3306-p079-b001,
  pdf-ffbe9c7b3306-p080-b006,
  pdf-ffbe9c7b3306-p081-b001,
  pdf-ffbe9c7b3306-p158-b001,
  pdf-ffbe9c7b3306-p158-b005,
  pdf-ffbe9c7b3306-p158-b006
}
\label{FABL.deterministicGoldreichLevinLearner}
\leanok
\uses{FABL.DeterministicGoldreichLevinInput, FABL.DeterministicQueryProgram, FABL.DeterministicQueryProgram.bind, FABL.Sign, FABL.SparseFourierHypothesis, FABL.deterministicGoldreichLevinProgram, FABL.deterministicGoldreichLevinSecondStage}
Given a deterministic Goldreich–Levin input, define a deterministic query program that
runs the Goldreich–Levin program and then, for each output state, runs the second-stage
program on that state, returning a sparse Fourier hypothesis in the input dimension.
\end{definition}

\begin{theorem}[Deterministic goldreich levin program concentrated family]
\lean{FABL.deterministicGoldreichLevinProgram_concentrated_family}
\statementsource{Analysis-of-Boolean-Functions-by-Ryan-ODonnell}{
  pdf-ffbe9c7b3306-p083-b002,
  pdf-ffbe9c7b3306-p083-b003
}
\label{FABL.deterministicGoldreichLevinProgram_concentrated_family}
\leanok
\uses{FABL.BooleanFunction, FABL.BooleanFunction.toReal, FABL.DeterministicGoldreichLevinInput, FABL.DeterministicGoldreichLevinInput.familySizeBound, FABL.DeterministicGoldreichLevinInput.familySizeBound_pos, FABL.DeterministicGoldreichLevinInput.learningParameter, FABL.DeterministicGoldreichLevinInput.threshold, FABL.DeterministicQueryProgram.runWithCost, FABL.IsFourierSpectrumConcentratedOn, FABL.card_l1ConcentratingFourierFamily_le_familySizeBound, FABL.deterministicGoldreichLevinProgram, FABL.deterministicGoldreichLevinProgram_isCorrectOutput, FABL.finiteFamily_nonempty_of_spectrum_concentrated, FABL.fourierCoeff, FABL.fourierOneNorm, FABL.goldreichLevinActiveCap, FABL.isFourierSpectrumConcentratedOn_l1ConcentratingFourierFamily, FABL.isFourierSpectrumConcentratedOn_of_goldreichLevin_complete, FABL.kushilevitzMansourThreshold, FABL.l1ConcentratingFourierFamily, FABL.positiveLearningParameter_toReal_mem_Ioc}
\difficulty{4}
% dataset-pipeline: stub -- prose pending build_informal_statements.py
Statement pending regeneration.
\end{theorem}

\begin{theorem}[Deterministic goldreich levin learner relative hamming dist le]
\lean{FABL.deterministicGoldreichLevinLearner_relativeHammingDist_le}
\statementsource{Analysis-of-Boolean-Functions-by-Ryan-ODonnell}{
  pdf-ffbe9c7b3306-p158-b006,
  pdf-ffbe9c7b3306-p080-b006,
  pdf-ffbe9c7b3306-p081-b001,
  pdf-ffbe9c7b3306-p080-b004,
  pdf-ffbe9c7b3306-p080-b005
}
\label{FABL.deterministicGoldreichLevinLearner_relativeHammingDist_le}
\leanok
\uses{FABL.BooleanFunction, FABL.BooleanFunction.toReal, FABL.DeterministicGoldreichLevinInput, FABL.DeterministicGoldreichLevinInput.coefficientInput, FABL.DeterministicGoldreichLevinInput.learningParameter, FABL.DeterministicQueryProgram.runWithCost, FABL.DeterministicQueryProgram.runWithCost_bind, FABL.DeterministicQueryProgram.runWithCost_deterministicGoldreichLevinCoefficientProgram, FABL.SmallBiasFourierInput.sample, FABL.SparseFourierHypothesis.evaluate, FABL.binaryCubeSignEquiv, FABL.deterministicGoldreichLevinHypothesis_coefficients_le, FABL.deterministicGoldreichLevinLearner, FABL.deterministicGoldreichLevinProgram, FABL.deterministicGoldreichLevinProgram_concentrated_family, FABL.deterministicGoldreichLevinSecondStage, FABL.fourierOneNorm, FABL.rationalSmallBiasFourierEstimate, FABL.relativeHammingDist, FABL.relativeHammingDist_sparseFourierHypothesis_of_coefficients_le}
\difficulty{4}
% dataset-pipeline: stub -- prose pending build_informal_statements.py
Statement pending regeneration.
\end{theorem}

\begin{definition}[Goldreich–Levin Learning Program]
\lean{FABL.deterministicGoldreichLevinLearningProgram}
\statementsource{Analysis-of-Boolean-Functions-by-Ryan-ODonnell}{
  pdf-ffbe9c7b3306-p079-b001,
  pdf-ffbe9c7b3306-p080-b006,
  pdf-ffbe9c7b3306-p081-b001,
  pdf-ffbe9c7b3306-p158-b006
}
\label{FABL.deterministicGoldreichLevinLearningProgram}
\leanok
\uses{FABL.DeterministicGoldreichLevinInput, FABL.DeterministicQueryProgram.toLearningProgram, FABL.LearningProgram, FABL.SparseFourierHypothesis, FABL.deterministicGoldreichLevinLearner}
Given a deterministic Goldreich–Levin input, define a membership-query learning program
in the input dimension by translating the deterministic Goldreich–Levin learner into a
learning program. Its output is a sparse Fourier hypothesis in that dimension.
\end{definition}

\begin{theorem}[Deterministic execution of the Goldreich–Levin learner]
\lean{FABL.runWithCost_deterministicGoldreichLevinLearningProgram}
\label{FABL.runWithCost_deterministicGoldreichLevinLearningProgram}
\leanok
\uses{FABL.BooleanFunction, FABL.DeterministicGoldreichLevinInput, FABL.DeterministicQueryProgram.runWithCost, FABL.DeterministicQueryProgram.runWithCost_toLearningProgram, FABL.LearningProgram.runWithCost, FABL.deterministicGoldreichLevinLearner, FABL.deterministicGoldreichLevinLearningProgram}
\difficulty{1}
Let $n\in\mathbb{N}$, let the learning input specify a dimension $n$, and let
$f:\{\pm1\}^n\to\{\pm1\}$ be a Boolean target. Running the deterministic Goldreich–Levin
learning program with target $f$ has the point-mass law at the output and cost returned
by executing its corresponding deterministic query program with oracle $f$.
\end{theorem}

\begin{theorem}[Deterministic goldreich levin learning program relative hamming dist le]
\lean{FABL.deterministicGoldreichLevinLearningProgram_relativeHammingDist_le}
\statementsource{Analysis-of-Boolean-Functions-by-Ryan-ODonnell}{
  pdf-ffbe9c7b3306-p080-b004,
  pdf-ffbe9c7b3306-p080-b005,
  pdf-ffbe9c7b3306-p080-b006,
  pdf-ffbe9c7b3306-p081-b001,
  pdf-ffbe9c7b3306-p158-b006
}
\label{FABL.deterministicGoldreichLevinLearningProgram_relativeHammingDist_le}
\leanok
\uses{FABL.BooleanFunction, FABL.BooleanFunction.toReal, FABL.DeterministicGoldreichLevinInput, FABL.DeterministicGoldreichLevinInput.learningParameter, FABL.LearningCost, FABL.LearningProgram.runWithCost, FABL.SparseFourierHypothesis, FABL.SparseFourierHypothesis.evaluate, FABL.deterministicGoldreichLevinLearner_relativeHammingDist_le, FABL.deterministicGoldreichLevinLearningProgram, FABL.fourierOneNorm, FABL.relativeHammingDist, FABL.runWithCost_deterministicGoldreichLevinLearningProgram}
\difficulty{3}
% dataset-pipeline: stub -- prose pending build_informal_statements.py
Statement pending regeneration.
\end{theorem}

\begin{theorem}[Cost of a Goldreich–Levin stage step]
\lean{FABL.DeterministicQueryProgram.runWithCost_deterministicGoldreichLevinStageStepProgram}
\label{FABL.DeterministicQueryProgram.runWithCost_deterministicGoldreichLevinStageStepProgram}
\leanok
\uses{FABL.BooleanFunction, FABL.DeterministicGoldreichLevinInput, FABL.DeterministicGoldreichLevinInput.threshold, FABL.DeterministicQueryProgram.map, FABL.DeterministicQueryProgram.runWithCost, FABL.DeterministicQueryProgram.runWithCost_bind, FABL.DeterministicQueryProgram.runWithCost_deterministicGoldreichLevinStageProgram, FABL.GoldreichLevinQueryState, FABL.deterministicGoldreichLevinStageCost, FABL.deterministicGoldreichLevinStageRecords, FABL.deterministicGoldreichLevinStageStepProgram, FABL.goldreichLevinControllerStepWork, FABL.goldreichLevinQueryStep}
\difficulty{3}
Let $n$ be a positive integer, let $a,b$ be positive integers with $2a\le b$, and let
$M\ge1$ be an integer Fourier norm bound. For a Boolean target
$f:\{\pm1\}^n\to\{\pm1\}$, a stage index $k<n$, a Goldreich–Levin query state, and a
finite active family of coordinate subsets, running the stage-step program returns the
state obtained by appending the stage’s estimate records and applying the prescribed
state update. Its cost is the sum of the stage cost and the controller’s charged work
for this dimension and threshold.
\end{theorem}

\begin{definition}[Deterministic stage query budget]
\lean{FABL.deterministicGoldreichLevinStageQueryBudget}
\label{FABL.deterministicGoldreichLevinStageQueryBudget}
\leanok
\uses{FABL.DeterministicGoldreichLevinInput, FABL.DeterministicGoldreichLevinInput.restrictedWeightInput, FABL.deterministicRestrictedWeightQueryCount}
Let the output dimension be a positive integer $n$, let the requested bias be a positive
rational number at most $1/2$, and let the integral Fourier $1$-norm bound be an integer
at least $1$. Define the stage query budget to be the restricted-weight query count.
\end{definition}

\begin{definition}[Deterministic stage work budget]
\lean{FABL.deterministicGoldreichLevinStageWorkBudget}
\label{FABL.deterministicGoldreichLevinStageWorkBudget}
\leanok
\uses{FABL.DeterministicGoldreichLevinInput, FABL.DeterministicGoldreichLevinInput.threshold, FABL.deterministicGoldreichLevinRestrictedLocalWork, FABL.deterministicGoldreichLevinStageQueryBudget, FABL.goldreichLevinActiveCap, FABL.goldreichLevinControllerStepWork}
Let the output dimension be a positive integer $n$, let the requested bias be a positive
rational number at most $1/2$, and let the integral Fourier $1$-norm bound be an integer
at least $1$. Define the stage work budget as the controller’s step work at the input
threshold, plus the stage query budget, plus twice the product of the active-set cap at
that threshold and the restricted-weight local work.
\end{definition}

\begin{definition}[Deterministic list query budget]
\lean{FABL.deterministicGoldreichLevinListQueryBudget}
\label{FABL.deterministicGoldreichLevinListQueryBudget}
\leanok
\uses{FABL.DeterministicGoldreichLevinInput, FABL.deterministicGoldreichLevinStageQueryBudget}
Let the output dimension be a positive integer $n$, let the requested bias be a positive
rational number at most $1/2$, and let the integral Fourier $1$-norm bound be an integer
at least $1$. Define the list query budget as $n$ times the stage query budget.
\end{definition}

\begin{definition}[Deterministic list work budget]
\lean{FABL.deterministicGoldreichLevinListWorkBudget}
\label{FABL.deterministicGoldreichLevinListWorkBudget}
\leanok
\uses{FABL.DeterministicGoldreichLevinInput, FABL.DeterministicGoldreichLevinInput.restrictedWeightInput, FABL.deterministicGoldreichLevinStageWorkBudget, FABL.deterministicRestrictedWeightConstructionWork}
Let the output dimension be a positive integer $n$, let the requested bias be a positive
rational number at most $1/2$, and let the integral Fourier $1$-norm bound be an integer
at least $1$. Define the list work budget as the restricted-weight construction work
plus $n$ times the stage work budget.
\end{definition}

\begin{definition}[Deterministic coefficient query budget]
\lean{FABL.deterministicGoldreichLevinCoefficientQueryBudget}
\label{FABL.deterministicGoldreichLevinCoefficientQueryBudget}
\leanok
\uses{FABL.DeterministicGoldreichLevinInput, FABL.DeterministicGoldreichLevinInput.coefficientInput, FABL.SmallBiasFourierInput.sampleCount}
Let the output dimension be a positive integer $n$, let the requested bias be a positive
rational number at most $1/2$, and let the integral Fourier $1$-norm bound be an integer
at least $1$. Define the coefficient query budget to be the coefficient estimator’s
sample count.
\end{definition}

\begin{definition}[Coefficient work budget]
\lean{FABL.deterministicGoldreichLevinCoefficientWorkBudget}
\label{FABL.deterministicGoldreichLevinCoefficientWorkBudget}
\leanok
\uses{FABL.DeterministicGoldreichLevinInput, FABL.DeterministicGoldreichLevinInput.coefficientInput, FABL.DeterministicGoldreichLevinInput.threshold, FABL.SmallBiasFourierInput.generatorInput, FABL.SmallBiasFourierInput.sampleCount, FABL.deterministicGoldreichLevinCoefficientLocalWork, FABL.deterministicSmallBiasWork, FABL.goldreichLevinActiveCap}
For a deterministic Goldreich–Levin input, define the coefficient work budget to be the
small-bias generator work plus the sample count, plus the active cap at the threshold
multiplied by the local work per coefficient.
\end{definition}

\begin{definition}[Query budget]
\lean{FABL.deterministicGoldreichLevinQueryBudget}
\label{FABL.deterministicGoldreichLevinQueryBudget}
\leanok
\uses{FABL.DeterministicGoldreichLevinInput, FABL.deterministicGoldreichLevinCoefficientQueryBudget, FABL.deterministicGoldreichLevinListQueryBudget}
For a deterministic Goldreich–Levin input, define the query budget as the sum of the
list query budget and the coefficient query budget.
\end{definition}

\begin{definition}[Work budget]
\lean{FABL.deterministicGoldreichLevinWorkBudget}
\label{FABL.deterministicGoldreichLevinWorkBudget}
\leanok
\uses{FABL.DeterministicGoldreichLevinInput, FABL.deterministicGoldreichLevinCoefficientWorkBudget, FABL.deterministicGoldreichLevinListWorkBudget}
For a deterministic Goldreich–Levin input, define the work budget as the sum of the list
work budget and the coefficient work budget.
\end{definition}

\begin{theorem}[Projected cost bound for the run]
\lean{FABL.deterministicGoldreichLevinRunUpToProgram_costProjection_le}
\label{FABL.deterministicGoldreichLevinRunUpToProgram_costProjection_le}
\leanok
\uses{FABL.BooleanFunction, FABL.DeterministicGoldreichLevinInput, FABL.DeterministicGoldreichLevinInput.threshold, FABL.DeterministicQueryProgram.runWithCost, FABL.DeterministicQueryProgram.runWithCost_bind, FABL.DeterministicQueryProgram.runWithCost_deterministicGoldreichLevinStageStepProgram, FABL.LearningCost, FABL.deterministicGoldreichLevinRunUpToProgram, FABL.deterministicGoldreichLevinStageCost, FABL.deterministicGoldreichLevinStageStepProgram, FABL.goldreichLevinActiveCap, FABL.goldreichLevinControllerStepWork}
\difficulty{5}
Let an input consist of a positive dimension $n$, positive integers $a,d$ with $2a\le
d$, and a Fourier bound $B\ge1$, and let the target be a function from the
$n$-dimensional sign cube to the signs. Suppose a map from learning costs to nonnegative
integers sends zero to zero and preserves addition. If, for some natural number $S$, the
projection of the controller cost plus the stage cost is at most $S$ at every stage with
index $k<n$ and active-set size at most the threshold's active cap, then for every $k\le
n$ the projected cost of running the program through $k$ stages is at most $kS$.
\end{theorem}

\begin{theorem}[Query bound for the run]
\lean{FABL.deterministicGoldreichLevinRunUpToProgram_queries_le}
\label{FABL.deterministicGoldreichLevinRunUpToProgram_queries_le}
\leanok
\uses{FABL.BooleanFunction, FABL.DeterministicGoldreichLevinInput, FABL.DeterministicGoldreichLevinInput.restrictedWeightInput, FABL.DeterministicQueryProgram.queryBatchCost, FABL.DeterministicQueryProgram.runWithCost, FABL.deterministicGoldreichLevinRunUpToProgram, FABL.deterministicGoldreichLevinRunUpToProgram_costProjection_le, FABL.deterministicGoldreichLevinStageCost, FABL.deterministicGoldreichLevinStageQueryBudget, FABL.deterministicRestrictedWeightQueryCount}
\difficulty{4}
Let an input consist of a positive dimension $n$, positive integers $a,d$ with $2a\le
d$, and a Fourier bound $B\ge1$, and let the target be a function from the
$n$-dimensional sign cube to the signs. For every natural number $k\le n$, the number of
membership queries made by running the program through $k$ stages is at most $k$ times
its stage query budget.
\end{theorem}

\begin{theorem}[First-Stage Work Bound]
\lean{FABL.deterministicGoldreichLevinRunUpToProgram_work_le}
\label{FABL.deterministicGoldreichLevinRunUpToProgram_work_le}
\leanok
\uses{FABL.BooleanFunction, FABL.DeterministicGoldreichLevinInput, FABL.DeterministicGoldreichLevinInput.threshold, FABL.DeterministicQueryProgram.runWithCost, FABL.card_goldreichLevinCandidates_le_two_mul, FABL.deterministicGoldreichLevinRestrictedLocalWork, FABL.deterministicGoldreichLevinRunUpToProgram, FABL.deterministicGoldreichLevinRunUpToProgram_costProjection_le, FABL.deterministicGoldreichLevinStageQueryBudget, FABL.deterministicGoldreichLevinStageWorkBudget, FABL.goldreichLevinCandidates, FABL.goldreichLevinControllerStepWork}
\difficulty{4}
Let the learning input have a positive dimension and a positive rational accuracy at
most one half, and let its integral Fourier norm bound be at least one. For every
Boolean function on the sign cube of that dimension and every integer $k$ no greater
than the dimension, the work used by the deterministic Goldreich–Levin program through
stage $k$ is at most $k$ times its per-stage work budget.
\end{theorem}

\begin{theorem}[No Random Examples in the First-Stage Program]
\lean{FABL.deterministicGoldreichLevinRunUpToProgram_randomExamples}
\label{FABL.deterministicGoldreichLevinRunUpToProgram_randomExamples}
\leanok
\uses{FABL.BooleanFunction, FABL.DeterministicGoldreichLevinInput, FABL.DeterministicQueryProgram.queryBatchCost, FABL.DeterministicQueryProgram.runWithCost, FABL.deterministicGoldreichLevinRunUpToProgram, FABL.deterministicGoldreichLevinRunUpToProgram_costProjection_le, FABL.deterministicGoldreichLevinStageCost}
\difficulty{3}
Let the learning input have a positive dimension and a positive rational accuracy at
most one half, and let its integral Fourier norm bound be at least one. For every
Boolean function on the sign cube of that dimension and every integer $k$ no greater
than the dimension, running the deterministic Goldreich–Levin program through stage $k$
uses zero random examples.
\end{theorem}

\begin{theorem}[Deterministic goldreich levin program resource bounds]
\lean{FABL.deterministicGoldreichLevinProgram_resource_bounds}
\statementsource{Analysis-of-Boolean-Functions-by-Ryan-ODonnell}{
  pdf-ffbe9c7b3306-p083-b002,
  pdf-ffbe9c7b3306-p085-b003
}
\label{FABL.deterministicGoldreichLevinProgram_resource_bounds}
\leanok
\uses{FABL.BooleanFunction, FABL.DeterministicGoldreichLevinInput, FABL.DeterministicGoldreichLevinInput.restrictedWeightInput, FABL.DeterministicQueryProgram.runWithCost, FABL.deterministicGoldreichLevinListQueryBudget, FABL.deterministicGoldreichLevinListWorkBudget, FABL.deterministicGoldreichLevinProgram, FABL.deterministicGoldreichLevinRunUpToProgram, FABL.deterministicGoldreichLevinRunUpToProgram_queries_le, FABL.deterministicGoldreichLevinRunUpToProgram_randomExamples, FABL.deterministicGoldreichLevinRunUpToProgram_work_le, FABL.deterministicRestrictedWeightConstructionWork}
\difficulty{4}
% dataset-pipeline: stub -- prose pending build_informal_statements.py
Statement pending regeneration.
\end{theorem}

\begin{theorem}[Coefficient program resource bounds]
\lean{FABL.deterministicGoldreichLevinCoefficientProgram_resource_bounds}
\label{FABL.deterministicGoldreichLevinCoefficientProgram_resource_bounds}
\leanok
\uses{FABL.BooleanFunction, FABL.DeterministicGoldreichLevinInput, FABL.DeterministicGoldreichLevinInput.coefficientInput, FABL.DeterministicGoldreichLevinInput.threshold, FABL.DeterministicQueryProgram.queryBatchCost, FABL.DeterministicQueryProgram.runWithCost, FABL.DeterministicQueryProgram.runWithCost_deterministicGoldreichLevinCoefficientProgram, FABL.SmallBiasFourierInput.generatorInput, FABL.SmallBiasFourierInput.sampleCount, FABL.deterministicGoldreichLevinCoefficientCost, FABL.deterministicGoldreichLevinCoefficientLocalWork, FABL.deterministicGoldreichLevinCoefficientProgram, FABL.deterministicGoldreichLevinCoefficientQueryBudget, FABL.deterministicGoldreichLevinCoefficientWorkBudget, FABL.deterministicSmallBiasWork, FABL.goldreichLevinActiveCap}
\difficulty{3}
Let a deterministic Goldreich–Levin input have dimension $n>0$, rational accuracy $a/b$
with $a>0$ and $2a\le b$, and Fourier bound $B\ge1$. For any Boolean target on $n$
sign-valued coordinates and any finite family of coordinate subsets whose size is at
most the input's active cap, the coefficient program makes zero random-example calls,
makes exactly the associated sample count of membership queries, and uses work at most
its coefficient work budget.
\end{theorem}

\begin{theorem}[Second-Stage Resource Bounds]
\lean{FABL.deterministicGoldreichLevinSecondStage_resource_bounds}
\label{FABL.deterministicGoldreichLevinSecondStage_resource_bounds}
\leanok
\uses{FABL.BooleanFunction, FABL.DeterministicGoldreichLevinInput, FABL.DeterministicGoldreichLevinInput.threshold, FABL.DeterministicQueryProgram.runWithCost, FABL.GoldreichLevinQueryState, FABL.deterministicGoldreichLevinCoefficientProgram_resource_bounds, FABL.deterministicGoldreichLevinCoefficientQueryBudget, FABL.deterministicGoldreichLevinCoefficientWorkBudget, FABL.deterministicGoldreichLevinSecondStage, FABL.goldreichLevinActiveCap}
\difficulty{4}
Let the learning input have a positive dimension and a positive rational accuracy at
most one half, and let its integral Fourier norm bound be at least one. For every
Boolean function on the sign cube of that dimension and every query state, the
second-stage program uses zero random examples, makes at most the coefficient sample
count in membership queries, and uses at most its coefficient work budget in local work.
\end{theorem}

\begin{theorem}[Learner resource bounds]
\lean{FABL.deterministicGoldreichLevinLearner_resource_bounds}
\label{FABL.deterministicGoldreichLevinLearner_resource_bounds}
\leanok
\uses{FABL.BooleanFunction, FABL.DeterministicGoldreichLevinInput, FABL.DeterministicQueryProgram.runWithCost, FABL.DeterministicQueryProgram.runWithCost_bind, FABL.deterministicGoldreichLevinLearner, FABL.deterministicGoldreichLevinProgram, FABL.deterministicGoldreichLevinProgram_resource_bounds, FABL.deterministicGoldreichLevinQueryBudget, FABL.deterministicGoldreichLevinSecondStage, FABL.deterministicGoldreichLevinSecondStage_resource_bounds, FABL.deterministicGoldreichLevinWorkBudget}
\difficulty{4}
Let $n$, $a$, $b$, and $s$ be natural numbers with $0<n$, $0<a$, $2a\le b$, and $1\le
s$. For every target function $f:\{\pm1\}^n\to\{\pm1\}$, the deterministic
Goldreich–Levin learner with dimension $n$, accuracy parameter $a/b$, and Fourier bound
$s$ uses zero random examples, makes at most its specified query budget, and performs at
most its specified work budget.
\end{theorem}

\begin{theorem}[Resource bounds for deterministic learning]
\lean{FABL.deterministicGoldreichLevinLearningProgram_resource_bounds}
\label{FABL.deterministicGoldreichLevinLearningProgram_resource_bounds}
\leanok
\uses{FABL.BooleanFunction, FABL.DeterministicGoldreichLevinInput, FABL.LearningCost, FABL.LearningProgram.runWithCost, FABL.SparseFourierHypothesis, FABL.deterministicGoldreichLevinLearner_resource_bounds, FABL.deterministicGoldreichLevinLearningProgram, FABL.deterministicGoldreichLevinQueryBudget, FABL.deterministicGoldreichLevinWorkBudget, FABL.runWithCost_deterministicGoldreichLevinLearningProgram}
\difficulty{3}
Let $n$ be the positive output dimension, and let $a,b$ be positive integers with $2a\le
b$. Let $M\ge 1$ be an integer Fourier $1$-norm bound, and let $f$ be a sign-valued
Boolean function on the $n$-dimensional cube. Every pair consisting of a sparse Fourier
hypothesis and a learning cost that has positive probability under the learner for $f$
uses zero random examples, at most the learner’s query budget, and at most the learner’s
work budget.
\end{theorem}

\begin{theorem}[Ceiling bound for the input scale]
\lean{FABL.smallBiasInput_scale_cast_le_dimension_div_epsilon_add_one}
\label{FABL.smallBiasInput_scale_cast_le_dimension_div_epsilon_add_one}
\leanok
\uses{FABL.SmallBiasInput, FABL.SmallBiasInput.epsilon, FABL.SmallBiasInput.scale}
\difficulty{5}
Let $n$, $a$, and $c$ be positive integers with $2a\leq c$. Define $\varepsilon=a/c$ and
$s=\lceil nc/a\rceil$. Then $s\leq n/\varepsilon+1$.
\end{theorem}

\begin{definition}[Runtime scale]
\lean{FABL.DeterministicGoldreichLevinInput.runtimeScale}
\label{FABL.DeterministicGoldreichLevinInput.runtimeScale}
\leanok
\uses{FABL.DeterministicGoldreichLevinInput, FABL.DeterministicGoldreichLevinInput.learningParameter}
For a deterministic Goldreich–Levin input with dimension $n$, Fourier bound $B$, and
learning parameter $\rho$, define its runtime scale to be $\frac{(n+1)(B+1)}{\rho}$.
\end{definition}

\begin{theorem}[Runtime scale factor bounds]
\lean{FABL.DeterministicGoldreichLevinInput.runtimeScale_factor_bounds}
\label{FABL.DeterministicGoldreichLevinInput.runtimeScale_factor_bounds}
\leanok
\uses{FABL.DeterministicGoldreichLevinInput, FABL.DeterministicGoldreichLevinInput.learningParameter, FABL.DeterministicGoldreichLevinInput.runtimeScale, FABL.positiveLearningParameter_toReal_mem_Ioc}
\difficulty{5}
Let $n,a,d,s$ be positive integers with $2a\le d$ and $s\ge1$, and set $\varepsilon=a/d$
and $R=(n+1)(s+1)/\varepsilon$. Then $R\ge1$, $n+1\le R$, $s+1\le R$, and
$1/\varepsilon\le R$.
\end{theorem}

\begin{theorem}[Threshold as a real quotient]
\lean{FABL.DeterministicGoldreichLevinInput.threshold_cast}
\label{FABL.DeterministicGoldreichLevinInput.threshold_cast}
\leanok
\uses{FABL.DeterministicGoldreichLevinInput, FABL.DeterministicGoldreichLevinInput.familySizeBound, FABL.DeterministicGoldreichLevinInput.learningParameter, FABL.DeterministicGoldreichLevinInput.threshold, FABL.kushilevitzMansourThreshold}
\difficulty{1}
Let $n,a,d,s$ be positive integers with $2a\le d$ and $s\ge1$. Let $\varepsilon=a/d$,
and let $M=\max\{1,\lceil4s^2/\varepsilon\rceil\}$. The threshold $\varepsilon/(4M)$,
viewed as a real number, equals $\varepsilon/(4M)$.
\end{theorem}

\begin{theorem}[Family size bound]
\lean{FABL.DeterministicGoldreichLevinInput.familySizeBound_add_one_cast_le_runtimeScale}
\label{FABL.DeterministicGoldreichLevinInput.familySizeBound_add_one_cast_le_runtimeScale}
\leanok
\uses{FABL.DeterministicGoldreichLevinInput, FABL.DeterministicGoldreichLevinInput.familySizeBound, FABL.DeterministicGoldreichLevinInput.learningParameter, FABL.DeterministicGoldreichLevinInput.runtimeScale, FABL.DeterministicGoldreichLevinInput.runtimeScale_factor_bounds, FABL.fourierOneNormFamilySizeBound_add_one_cast_le, FABL.positiveLearningParameter_toReal_mem_Ioc}
\difficulty{5}
Let a deterministic Goldreich–Levin input have dimension $n>0$, accuracy $a/b$ with
integers $a>0$ and $2a\le b$, and an integer Fourier bound $S\ge1$. Set
$\varepsilon=a/b$, $M=\max\!\left(1,\left\lceil 4S^2/\varepsilon\right\rceil\right)$,
and $R=(n+1)(S+1)/\varepsilon$. Then $M+1\le 6R^3$.
\end{theorem}

\begin{theorem}[Active cap bound by runtime scale]
\lean{FABL.DeterministicGoldreichLevinInput.activeCap_add_one_cast_le_runtimeScale}
\label{FABL.DeterministicGoldreichLevinInput.activeCap_add_one_cast_le_runtimeScale}
\leanok
\uses{FABL.DeterministicGoldreichLevinInput, FABL.DeterministicGoldreichLevinInput.familySizeBound, FABL.DeterministicGoldreichLevinInput.familySizeBound_add_one_cast_le_runtimeScale, FABL.DeterministicGoldreichLevinInput.learningParameter, FABL.DeterministicGoldreichLevinInput.runtimeScale, FABL.DeterministicGoldreichLevinInput.runtimeScale_factor_bounds, FABL.DeterministicGoldreichLevinInput.threshold, FABL.DeterministicGoldreichLevinInput.threshold_cast, FABL.goldreichLevinActiveCap, FABL.goldreichLevinActiveCap_add_one_cast_le, FABL.positiveLearningParameter_toReal_mem_Ioc}
\difficulty{5}
Let $n,a,d,B$ be integers with $n>0$, $a>0$, $2a\le d$, and $B\ge1$. Set
$\varepsilon=a/d$, $M=\max\{1,\lceil 4B^2/\varepsilon\rceil\}$, $\tau=\varepsilon/(4M)$,
and $R=(n+1)(B+1)/\varepsilon$. Then $\lceil4/\tau^2\rceil+1\le 2^{13}R^8$.
\end{theorem}

\begin{theorem}[Restricted-component polynomial bounds]
\lean{FABL.deterministicGoldreichLevinRestrictedComponents_le_polynomialBudget}
\label{FABL.deterministicGoldreichLevinRestrictedComponents_le_polynomialBudget}
\leanok
\uses{FABL.DeterministicGoldreichLevinInput, FABL.DeterministicGoldreichLevinInput.restrictedWeightInput, FABL.RestrictedWeightInput.polynomialBudget, FABL.deterministicGoldreichLevinRestrictedLocalWork, FABL.deterministicGoldreichLevinStageQueryBudget, FABL.deterministicRestrictedWeightConstructionWork, FABL.deterministicRestrictedWeightCost_queries, FABL.deterministicRestrictedWeightCost_resource_bounds, FABL.deterministicRestrictedWeightCost_work, FABL.deterministicRestrictedWeightLocalWork, FABL.deterministicRestrictedWeightQueryCount}
\difficulty{4}
Let an input consist of a positive dimension $n$, positive integers $a,d$ with $2a\le
d$, and an integer Fourier bound $B\ge1$. For its associated restricted-weight input,
let $A$ be the sum of the inner and outer algorithm scales; its polynomial budget is
$2^{20}(A+1)^8$. The restricted-weight construction work, the stage query budget, and
the restricted local work are each at most this budget.
\end{theorem}

\begin{theorem}[Coefficient components fit a polynomial budget]
\lean{FABL.deterministicGoldreichLevinCoefficientComponents_le_polynomialBudget}
\label{FABL.deterministicGoldreichLevinCoefficientComponents_le_polynomialBudget}
\leanok
\uses{FABL.DeterministicGoldreichLevinInput, FABL.DeterministicGoldreichLevinInput.coefficientInput, FABL.SmallBiasFourierInput.generatorInput, FABL.SmallBiasFourierInput.polynomialBudget, FABL.SmallBiasFourierInput.sampleCount, FABL.deterministicGoldreichLevinCoefficientLocalWork, FABL.deterministicSmallBiasFourierEstimatorCost_queries, FABL.deterministicSmallBiasFourierEstimatorCost_resource_bounds, FABL.deterministicSmallBiasFourierEstimatorCost_work, FABL.deterministicSmallBiasWork, FABL.smallBiasFourierEstimatorWork}
\difficulty{3}
Let a deterministic Goldreich–Levin input have dimension $n>0$, rational accuracy $a/b$
with $a>0$ and $2a\le b$, and Fourier bound $B\ge1$. Let $P$ be the polynomial budget
associated with its coefficient input. The deterministic small-bias construction work at
the associated generator parameters, the associated sample count, and the local
Fourier-estimation work over all $n$ coordinates are each at most $P$.
\end{theorem}

\begin{theorem}[Restricted weight scale bound]
\lean{FABL.DeterministicGoldreichLevinInput.restrictedWeightAlgorithmScale_add_one_cast_le_runtimeScale}
\label{FABL.DeterministicGoldreichLevinInput.restrictedWeightAlgorithmScale_add_one_cast_le_runtimeScale}
\leanok
\uses{FABL.DeterministicGoldreichLevinInput, FABL.DeterministicGoldreichLevinInput.familySizeBound, FABL.DeterministicGoldreichLevinInput.familySizeBound_add_one_cast_le_runtimeScale, FABL.DeterministicGoldreichLevinInput.familySizeBound_pos, FABL.DeterministicGoldreichLevinInput.learningParameter, FABL.DeterministicGoldreichLevinInput.restrictedWeightInput, FABL.DeterministicGoldreichLevinInput.restrictedWeightInput_epsilon, FABL.DeterministicGoldreichLevinInput.runtimeScale, FABL.DeterministicGoldreichLevinInput.runtimeScale_factor_bounds, FABL.DeterministicGoldreichLevinInput.threshold, FABL.DeterministicGoldreichLevinInput.threshold_cast, FABL.RestrictedWeightInput.algorithmScale, FABL.RestrictedWeightInput.epsilon, FABL.RestrictedWeightInput.innerInput, FABL.RestrictedWeightInput.innerInput_biasParameter, FABL.RestrictedWeightInput.outerInput, FABL.RestrictedWeightInput.outerInput_biasParameter, FABL.SmallBiasFourierInput.epsilon, FABL.SmallBiasFourierInput.generatorInput, FABL.SmallBiasFourierInput.generatorInput_epsilon, FABL.SmallBiasInput.epsilon, FABL.SmallBiasInput.scale, FABL.positiveLearningParameter_toReal_mem_Ioc, FABL.smallBiasInput_scale_cast_le_dimension_div_epsilon_add_one}
\difficulty{6}
Let a deterministic Goldreich–Levin input have dimension $n>0$, accuracy $a/b$ with
integers $a>0$ and $2a\le b$, and an integer Fourier bound $S\ge1$. Put
$\varepsilon=a/b$ and $R=(n+1)(S+1)/\varepsilon$. The sum $A=\left\lceil
4nbS/a\right\rceil+\left\lceil 4nbS^2/a\right\rceil$ of the inner and outer small-bias
scales satisfies $A+1\le 2^{16}R^{11}$.
\end{theorem}

\begin{theorem}[Restricted weight budget bound]
\lean{FABL.DeterministicGoldreichLevinInput.restrictedWeightPolynomialBudget_cast_le_runtimeScale}
\label{FABL.DeterministicGoldreichLevinInput.restrictedWeightPolynomialBudget_cast_le_runtimeScale}
\leanok
\uses{FABL.DeterministicGoldreichLevinInput, FABL.DeterministicGoldreichLevinInput.restrictedWeightAlgorithmScale_add_one_cast_le_runtimeScale, FABL.DeterministicGoldreichLevinInput.restrictedWeightInput, FABL.DeterministicGoldreichLevinInput.runtimeScale, FABL.RestrictedWeightInput.algorithmScale, FABL.RestrictedWeightInput.polynomialBudget}
\difficulty{3}
Let $n,a,d,s$ be positive integers with $2a\le d$ and $s\ge1$, and put
$\varepsilon=a/d$. Define $M=\left\lceil n\,4d s/a\right\rceil+\left\lceil n\,4d
s^2/a\right\rceil$ and $R=(n+1)(s+1)/\varepsilon$. Then $2^{20}(M+1)^8\le
2^{148}R^{88}$.
\end{theorem}

\begin{theorem}[Coefficient generator scale bound]
\lean{FABL.DeterministicGoldreichLevinInput.coefficientGeneratorScale_add_one_cast_le_runtimeScale}
\label{FABL.DeterministicGoldreichLevinInput.coefficientGeneratorScale_add_one_cast_le_runtimeScale}
\leanok
\uses{FABL.DeterministicGoldreichLevinInput, FABL.DeterministicGoldreichLevinInput.activeCap_add_one_cast_le_runtimeScale, FABL.DeterministicGoldreichLevinInput.coefficientInput, FABL.DeterministicGoldreichLevinInput.coefficientInput_epsilon, FABL.DeterministicGoldreichLevinInput.learningParameter, FABL.DeterministicGoldreichLevinInput.runtimeScale, FABL.DeterministicGoldreichLevinInput.runtimeScale_factor_bounds, FABL.DeterministicGoldreichLevinInput.threshold, FABL.SmallBiasFourierInput.epsilon, FABL.SmallBiasFourierInput.generatorInput, FABL.SmallBiasFourierInput.generatorInput_epsilon, FABL.SmallBiasInput.epsilon, FABL.SmallBiasInput.scale, FABL.goldreichLevinActiveCap, FABL.goldreichLevinActiveCap_pos, FABL.positiveLearningParameter_toReal_mem_Ioc, FABL.smallBiasInput_scale_cast_le_dimension_div_epsilon_add_one}
\difficulty{5}
For every deterministic Goldreich–Levin input, let $S$ denote the scale of the generator
input associated to its coefficient input, and let $R$ be its runtime scale. Then
$S+1\le 2^{15}R^{11}$.
\end{theorem}

\begin{theorem}[Polynomial budget bound]
\lean{FABL.DeterministicGoldreichLevinInput.coefficientPolynomialBudget_cast_le_runtimeScale}
\label{FABL.DeterministicGoldreichLevinInput.coefficientPolynomialBudget_cast_le_runtimeScale}
\leanok
\uses{FABL.DeterministicGoldreichLevinInput, FABL.DeterministicGoldreichLevinInput.coefficientGeneratorScale_add_one_cast_le_runtimeScale, FABL.DeterministicGoldreichLevinInput.coefficientInput, FABL.DeterministicGoldreichLevinInput.runtimeScale, FABL.SmallBiasFourierInput.generatorInput, FABL.SmallBiasFourierInput.polynomialBudget, FABL.SmallBiasInput.scale}
\difficulty{4}
Let a deterministic Goldreich–Levin input have dimension $n>0$, accuracy $a/b$ with
integers $a>0$ and $2a\le b$, and an integer Fourier bound $S\ge1$. Put
$\varepsilon=a/b$, $M=\max\!\left(1,\left\lceil 4S^2/\varepsilon\right\rceil\right)$,
and $R=(n+1)(S+1)/\varepsilon$. For $d=\left\lceil n(4bS)/a\right\rceil$, the
coefficient input’s polynomial budget $2^{18}(d+1)^8$ is at most $2^{138}R^{88}$.
\end{theorem}

\begin{theorem}[Restricted-component runtime bounds]
\lean{FABL.deterministicGoldreichLevinRestrictedComponents_cast_le_runtimeScale}
\label{FABL.deterministicGoldreichLevinRestrictedComponents_cast_le_runtimeScale}
\leanok
\uses{FABL.DeterministicGoldreichLevinInput, FABL.DeterministicGoldreichLevinInput.restrictedWeightInput, FABL.DeterministicGoldreichLevinInput.restrictedWeightPolynomialBudget_cast_le_runtimeScale, FABL.DeterministicGoldreichLevinInput.runtimeScale, FABL.RestrictedWeightInput.polynomialBudget, FABL.deterministicGoldreichLevinRestrictedComponents_le_polynomialBudget, FABL.deterministicGoldreichLevinRestrictedLocalWork, FABL.deterministicGoldreichLevinStageQueryBudget, FABL.deterministicRestrictedWeightConstructionWork}
\difficulty{4}
Let an input consist of a positive dimension $n$, positive integers $a,d$ with $2a\le
d$, and an integer Fourier bound $B\ge1$. Put $\varepsilon=a/d$ and
$R=(n+1)(B+1)/\varepsilon$. The restricted-weight construction work, the stage query
budget, and the restricted local work are each at most $2^{148}R^{88}$ when viewed as
real numbers.
\end{theorem}

\begin{theorem}[Coefficient routine runtime bounds]
\lean{FABL.deterministicGoldreichLevinCoefficientComponents_cast_le_runtimeScale}
\label{FABL.deterministicGoldreichLevinCoefficientComponents_cast_le_runtimeScale}
\leanok
\uses{FABL.DeterministicGoldreichLevinInput, FABL.DeterministicGoldreichLevinInput.coefficientInput, FABL.DeterministicGoldreichLevinInput.coefficientPolynomialBudget_cast_le_runtimeScale, FABL.DeterministicGoldreichLevinInput.runtimeScale, FABL.SmallBiasFourierInput.generatorInput, FABL.SmallBiasFourierInput.polynomialBudget, FABL.SmallBiasFourierInput.sampleCount, FABL.deterministicGoldreichLevinCoefficientComponents_le_polynomialBudget, FABL.deterministicGoldreichLevinCoefficientLocalWork, FABL.deterministicSmallBiasWork}
\difficulty{4}
Let a deterministic Goldreich–Levin input have dimension $n>0$, rational accuracy $a/b$
with $a>0$ and $2a\le b$, and Fourier bound $B\ge1$. Write $R$ for its runtime scale,
$W$ for the deterministic small-bias construction work at the associated generator
parameters, $N$ for the associated sample count, and $L$ for the local
Fourier-estimation work over all $n$ coordinates. Then $W\le 2^{138}R^{88}$, $N\le
2^{138}R^{88}$, and $L\le 2^{138}R^{88}$.
\end{theorem}

\begin{definition}[Polynomial Runtime Bound]
\lean{FABL.DeterministicGoldreichLevinInput.polynomialRuntimeBound}
\statementsource{Analysis-of-Boolean-Functions-by-Ryan-ODonnell}{
  pdf-ffbe9c7b3306-p083-b002,
  pdf-ffbe9c7b3306-p085-b003,
  pdf-ffbe9c7b3306-p158-b006
}
\label{FABL.DeterministicGoldreichLevinInput.polynomialRuntimeBound}
\leanok
\uses{FABL.DeterministicGoldreichLevinInput, FABL.DeterministicGoldreichLevinInput.runtimeScale}
For a deterministic Goldreich–Levin input, define its polynomial runtime bound to be
$2^{170}$ times the one-hundredth power of its runtime scale.
\end{definition}

\begin{theorem}[Query budget runtime bound]
\lean{FABL.DeterministicGoldreichLevinInput.queryBudget_cast_le_polynomialRuntimeBound}
\label{FABL.DeterministicGoldreichLevinInput.queryBudget_cast_le_polynomialRuntimeBound}
\leanok
\uses{FABL.DeterministicGoldreichLevinInput, FABL.DeterministicGoldreichLevinInput.coefficientInput, FABL.DeterministicGoldreichLevinInput.polynomialRuntimeBound, FABL.DeterministicGoldreichLevinInput.runtimeScale, FABL.DeterministicGoldreichLevinInput.runtimeScale_factor_bounds, FABL.SmallBiasFourierInput.sampleCount, FABL.deterministicGoldreichLevinCoefficientComponents_cast_le_runtimeScale, FABL.deterministicGoldreichLevinCoefficientQueryBudget, FABL.deterministicGoldreichLevinListQueryBudget, FABL.deterministicGoldreichLevinQueryBudget, FABL.deterministicGoldreichLevinRestrictedComponents_cast_le_runtimeScale, FABL.deterministicGoldreichLevinStageQueryBudget}
\difficulty{5}
Let a deterministic Goldreich–Levin input have dimension $n>0$, accuracy $a/b$ with
integers $a>0$ and $2a\le b$, and an integer Fourier bound $S\ge1$. Put
$\varepsilon=a/b$, $M=\max\!\left(1,\left\lceil 4S^2/\varepsilon\right\rceil\right)$,
$R=(n+1)(S+1)/\varepsilon$, and $D=2^{\max(1,\lceil\log_2(\lceil 4nb/a\rceil)\rceil)}$.
The total query budget is at most $2^{170}R^{100}$, where the budget is $n$ times the
product of the sample counts $2^{\deg_i}2^{\deg_o}$, plus the coefficient sample count
$2^{\deg_c}2^{\deg_c}$; here $\deg_i,\deg_o,\deg_c$ are respectively
$\max(1,\lceil\log_2\lceil 4nbS/a\rceil\rceil)$, $\max(1,\lceil\log_2\lceil
4nbS^2/a\rceil\rceil)$, and $\max(1,\lceil\log_2\lceil
2nb(2\lceil4/(\varepsilon/(4M))^2\rceil)/a\rceil\rceil)$.
\end{theorem}

\begin{theorem}[Work budget polynomial bound]
\lean{FABL.DeterministicGoldreichLevinInput.workBudget_cast_le_polynomialRuntimeBound}
\label{FABL.DeterministicGoldreichLevinInput.workBudget_cast_le_polynomialRuntimeBound}
\leanok
\uses{FABL.DeterministicGoldreichLevinInput, FABL.DeterministicGoldreichLevinInput.activeCap_add_one_cast_le_runtimeScale, FABL.DeterministicGoldreichLevinInput.coefficientInput, FABL.DeterministicGoldreichLevinInput.polynomialRuntimeBound, FABL.DeterministicGoldreichLevinInput.restrictedWeightInput, FABL.DeterministicGoldreichLevinInput.runtimeScale, FABL.DeterministicGoldreichLevinInput.runtimeScale_factor_bounds, FABL.DeterministicGoldreichLevinInput.threshold, FABL.SmallBiasFourierInput.generatorInput, FABL.SmallBiasFourierInput.sampleCount, FABL.deterministicGoldreichLevinCoefficientComponents_cast_le_runtimeScale, FABL.deterministicGoldreichLevinCoefficientLocalWork, FABL.deterministicGoldreichLevinCoefficientWorkBudget, FABL.deterministicGoldreichLevinListWorkBudget, FABL.deterministicGoldreichLevinRestrictedComponents_cast_le_runtimeScale, FABL.deterministicGoldreichLevinRestrictedLocalWork, FABL.deterministicGoldreichLevinStageQueryBudget, FABL.deterministicGoldreichLevinStageWorkBudget, FABL.deterministicGoldreichLevinWorkBudget, FABL.deterministicRestrictedWeightConstructionWork, FABL.deterministicSmallBiasWork, FABL.goldreichLevinActiveCap, FABL.goldreichLevinControllerStepWork}
\difficulty{6}
Let $n,a,d,s$ be positive integers with $2a\le d$ and $s\ge1$, set $\varepsilon=a/d$,
and define $R=(n+1)(s+1)/\varepsilon$. The deterministic Goldreich–Levin work budget for
these parameters is at most $2^{170}R^{100}$.
\end{theorem}

\begin{theorem}[Polynomial query bound]
\lean{FABL.deterministicGoldreichLevinLearner_queries_polynomial_le}
\proofsource{Analysis-of-Boolean-Functions-by-Ryan-ODonnell}{
  pdf-ffbe9c7b3306-p158-b006,
  pdf-ffbe9c7b3306-p080-b003,
  pdf-ffbe9c7b3306-p080-b006,
  pdf-ffbe9c7b3306-p081-b001
}
\label{FABL.deterministicGoldreichLevinLearner_queries_polynomial_le}
\leanok
\uses{FABL.BooleanFunction, FABL.DeterministicGoldreichLevinInput, FABL.DeterministicGoldreichLevinInput.polynomialRuntimeBound, FABL.DeterministicGoldreichLevinInput.queryBudget_cast_le_polynomialRuntimeBound, FABL.DeterministicQueryProgram.runWithCost, FABL.deterministicGoldreichLevinLearner, FABL.deterministicGoldreichLevinLearner_resource_bounds, FABL.deterministicGoldreichLevinQueryBudget}
\difficulty{3}
Let $n$, $a$, $b$, and $s$ be natural numbers with $0<n$, $0<a$, $2a\le b$, and $1\le
s$. For every target function $f:\{\pm1\}^n\to\{\pm1\}$, the deterministic
Goldreich–Levin learner with dimension $n$, accuracy parameter $a/b$, and Fourier bound
$s$ uses at most $2^{170}\bigl((n+1)(s+1)/(a/b)\bigr)^{100}$ membership queries.
\end{theorem}

\begin{theorem}[Polynomial work bound for the learner]
\lean{FABL.deterministicGoldreichLevinLearner_work_polynomial_le}
\label{FABL.deterministicGoldreichLevinLearner_work_polynomial_le}
\leanok
\uses{FABL.BooleanFunction, FABL.DeterministicGoldreichLevinInput, FABL.DeterministicGoldreichLevinInput.polynomialRuntimeBound, FABL.DeterministicGoldreichLevinInput.workBudget_cast_le_polynomialRuntimeBound, FABL.DeterministicQueryProgram.runWithCost, FABL.deterministicGoldreichLevinLearner, FABL.deterministicGoldreichLevinLearner_resource_bounds, FABL.deterministicGoldreichLevinWorkBudget}
\difficulty{3}
Let the dimension be $n>0$, let the accuracy be $\rho=a/b$ for positive integers $a,b$
with $2a\le b$, and let $M\ge1$ be an integer Fourier $1$-norm bound. For every
sign-valued Boolean target on the $n$-dimensional cube, the learner's charged work,
viewed as a real number, is at most $2^{170}\left((n+1)(M+1)/\rho\right)^{100}$.
\end{theorem}

\begin{theorem}[Deterministic goldreich levin learning program spec]
\lean{FABL.deterministicGoldreichLevinLearningProgram_spec}
\statementsource{Analysis-of-Boolean-Functions-by-Ryan-ODonnell}{
  pdf-ffbe9c7b3306-p158-b006,
  pdf-ffbe9c7b3306-p080-b006,
  pdf-ffbe9c7b3306-p081-b001
}
\label{FABL.deterministicGoldreichLevinLearningProgram_spec}
\leanok
\uses{FABL.BooleanFunction, FABL.BooleanFunction.toReal, FABL.DeterministicGoldreichLevinInput, FABL.DeterministicGoldreichLevinInput.learningParameter, FABL.DeterministicGoldreichLevinInput.polynomialRuntimeBound, FABL.DeterministicGoldreichLevinInput.queryBudget_cast_le_polynomialRuntimeBound, FABL.DeterministicGoldreichLevinInput.workBudget_cast_le_polynomialRuntimeBound, FABL.LearningCost, FABL.LearningProgram.runWithCost, FABL.SparseFourierHypothesis, FABL.SparseFourierHypothesis.evaluate, FABL.deterministicGoldreichLevinLearningProgram, FABL.deterministicGoldreichLevinLearningProgram_relativeHammingDist_le, FABL.deterministicGoldreichLevinLearningProgram_resource_bounds, FABL.deterministicGoldreichLevinQueryBudget, FABL.deterministicGoldreichLevinWorkBudget, FABL.fourierOneNorm, FABL.relativeHammingDist}
\difficulty{4}
% dataset-pipeline: stub -- prose pending build_informal_statements.py
Statement pending regeneration.
\end{theorem}

\begin{definition}[Deterministic Goldreich–Levin task]
\lean{FABL.DeterministicGoldreichLevinTask}
\statementsource{Analysis-of-Boolean-Functions-by-Ryan-ODonnell}{
  pdf-ffbe9c7b3306-p083-b002,
  pdf-ffbe9c7b3306-p083-b003
}
\label{FABL.DeterministicGoldreichLevinTask}
\leanok
\uses{FABL.BooleanFunction, FABL.DeterministicGoldreichLevinInput}
A deterministic Goldreich–Levin task consists of a deterministic Goldreich–Levin input
together with a Boolean function on the cube whose dimension is the input’s accuracy
dimension.
\end{definition}

\begin{definition}[Task scale]
\lean{FABL.deterministicGoldreichLevinTaskScale}
\label{FABL.deterministicGoldreichLevinTaskScale}
\leanok
\uses{FABL.DeterministicGoldreichLevinInput.runtimeScale, FABL.DeterministicGoldreichLevinTask}
A deterministic Goldreich–Levin task consists of a deterministic Goldreich–Levin input
together with a Boolean function on the sign cube of the input dimension. Define the
task scale to be the runtime scale of its input.
\end{definition}

\begin{theorem}[Deterministic goldreich levin learner queries is big o]
\lean{FABL.deterministicGoldreichLevinLearner_queries_isBigO}
\statementsource{Analysis-of-Boolean-Functions-by-Ryan-ODonnell}{pdf-ffbe9c7b3306-p158-b006}
\label{FABL.deterministicGoldreichLevinLearner_queries_isBigO}
\leanok
\uses{FABL.DeterministicGoldreichLevinInput.polynomialRuntimeBound, FABL.DeterministicGoldreichLevinInput.runtimeScale_factor_bounds, FABL.DeterministicGoldreichLevinTask, FABL.DeterministicQueryProgram.runWithCost, FABL.deterministicGoldreichLevinLearner, FABL.deterministicGoldreichLevinLearner_queries_polynomial_le, FABL.deterministicGoldreichLevinTaskScale}
\difficulty{4}
% dataset-pipeline: stub -- prose pending build_informal_statements.py
Statement pending regeneration.
\end{theorem}

\begin{theorem}[Deterministic goldreich levin learner work is big o]
\lean{FABL.deterministicGoldreichLevinLearner_work_isBigO}
\statementsource{Analysis-of-Boolean-Functions-by-Ryan-ODonnell}{
  pdf-ffbe9c7b3306-p080-b006,
  pdf-ffbe9c7b3306-p081-b001,
  pdf-ffbe9c7b3306-p158-b006
}
\label{FABL.deterministicGoldreichLevinLearner_work_isBigO}
\leanok
\uses{FABL.DeterministicGoldreichLevinInput.polynomialRuntimeBound, FABL.DeterministicGoldreichLevinInput.runtimeScale_factor_bounds, FABL.DeterministicGoldreichLevinTask, FABL.DeterministicQueryProgram.runWithCost, FABL.deterministicGoldreichLevinLearner, FABL.deterministicGoldreichLevinLearner_work_polynomial_le, FABL.deterministicGoldreichLevinTaskScale}
\difficulty{4}
% dataset-pipeline: stub -- prose pending build_informal_statements.py
Statement pending regeneration.
\end{theorem}
